\documentclass[11pt,a4paper]{article}

% ── Packages ──────────────────────────────────────────────────────────────────
\usepackage[utf8]{inputenc}
\usepackage[T1]{fontenc}
\usepackage{amsmath,amssymb,amsthm}
\usepackage{mathtools}
\usepackage{enumitem}
\usepackage[margin=2.5cm]{geometry}
\usepackage{hyperref}
\usepackage{cleveref}
\hypersetup{hypertexnames=false}
\usepackage{algorithm}
\usepackage{algpseudocode}

% ── Theorem environments ─────────────────────────────────────────────────────
\theoremstyle{plain}
\newtheorem{theorem}{Theorem}[section]
\newtheorem{proposition}[theorem]{Proposition}
\newtheorem{lemma}[theorem]{Lemma}
\newtheorem{corollary}[theorem]{Corollary}

\theoremstyle{definition}
\newtheorem{definition}[theorem]{Definition}
\newtheorem{example}[theorem]{Example}
\newtheorem{remark}[theorem]{Remark}

% ── Notation ─────────────────────────────────────────────────────────────────
\newcommand{\Nfwd}[1]{N^{+}[#1]}
\newcommand{\Nbwd}[1]{N^{-}[#1]}
\newcommand{\Ncl}[1]{N[#1]}
\newcommand{\taufwd}{\tau^{+}}
\newcommand{\taubwd}{\tau^{-}}
\newcommand{\powerset}{\mathcal{P}}
\newcommand{\subbase}{\mathcal{S}}
\newcommand{\base}{\mathcal{B}}

% ── Title ────────────────────────────────────────────────────────────────────
\title{%
  Bitopological Spaces from Directed Graphs:\\
  Extending the Nada Construction to Capture\\
  Temporal Irreversibility%
}

\author{%
  Jos\'e Mauricio G\'omez Juli\'an\\
  \textit{Independent Researcher}\\
  \texttt{mauriciogomezjulian@protonmail.ch}%
}

\date{\today}

% ══════════════════════════════════════════════════════════════════════════════
\begin{document}
\maketitle

\begin{abstract}
Nada et al.\ (2018) constructed a topology on the vertex set of an
undirected graph by taking the closed neighborhoods as a subbase and
closing under finite intersections and arbitrary unions.  We extend this
construction to directed graphs: for a digraph $G = (V, E)$, applying the
procedure separately to the forward and backward closed neighborhoods
yields two topologies $\taufwd$ and~$\taubwd$ on~$V$.  The triple
$(V, \taufwd, \taubwd)$ is a bitopological space in the sense of
Kelly~(1963); no hypothesis on~$G$ is required, because the topology
axioms are universal operations on the powerset.  The asymmetry between
$\taufwd$ and $\taubwd$ provides a topological quantification of temporal
irreversibility when~$G$ arises as the directed visibility graph of a
time series.  We prove the central theorem and related results on the
Alexandrov sublayer, the specialization preorder, and pairwise
connectedness; supply algorithms with polynomial-time complexity and a
three-valued pairwise decision procedure that honestly reports resource
exhaustion; formalize the core proofs in Lean~4 against Mathlib; and
apply the framework to quarterly U.S.\ real GDP growth (1992--2024,
$n = 129$), recovering a positive asymmetry direction consistent with
slow-expansion / abrupt-contraction dynamics and a coherent topological
classification of the COVID-19 contraction and rebound as a single
episode.
\end{abstract}

\medskip
\noindent\textbf{Keywords:} bitopological spaces, directed graphs, visibility
graphs, time series analysis, temporal irreversibility, Lean 4 formalization.

\medskip
\noindent\textbf{MSC 2020:} 54E55, 05C20, 37M10, 68V20.

% ══════════════════════════════════════════════════════════════════════════════
\section{Introduction}\label{sec:intro}
% ══════════════════════════════════════════════════════════════════════════════

The analysis of time series through network-theoretic methods has become a
fruitful line of research since the introduction of visibility graphs by Lacasa
et al.~\cite{Lacasa2008}.  Given a time series $(x_1, x_2, \ldots, x_n)$, one
constructs a graph whose vertices are the time indices and whose edges encode a
geometric visibility condition between data points.  Topological properties of
the resulting graph---degree distributions, clustering coefficients, community
structure---then serve as invariants of the underlying dynamics.

A natural refinement is to orient the edges of the visibility graph from earlier
to later time indices, producing a directed acyclic graph (DAG) in which the
time index serves as a topological order.  This orientation carries information
about temporal asymmetry that is invisible to any undirected construction: in
systems with asymmetric dynamics---sharp contractions followed by slow
recoveries in financial markets, rapid heating followed by gradual relaxation in
physical systems, causal precedence in event sequences---the directed graph
distinguishes the forward flow of time from the backward flow, while its
undirected shadow does not.

Nada et al.~\cite{Nada2018} introduced a procedure for endowing the vertex set
of an undirected graph with a topology: one takes the family of closed
neighborhoods $\{\Ncl{v} : v \in V\}$ as a subbase and generates a topology by
closing under finite intersections and arbitrary unions.  This construction
yields topological invariants (connectedness, separation, component counts) that
capture structural properties of the graph.

The central observation of this paper is that the Nada procedure extends without
modification to directed graphs and, in doing so, produces not one but two
topologies.  Given a digraph $G = (V, E)$, define:
\begin{itemize}[nosep]
  \item the \emph{forward closed neighborhood}
        $\Nfwd{v} = \{v\} \cup \{w : (v,w) \in E\}$;
  \item the \emph{backward closed neighborhood}
        $\Nbwd{v} = \{v\} \cup \{w : (w,v) \in E\}$.
\end{itemize}
Applying the Nada procedure to $\subbase^{+} = \{\Nfwd{v} : v \in V\}$ yields a
topology $\taufwd$; applying it to $\subbase^{-} = \{\Nbwd{v} : v \in V\}$
yields a topology $\taubwd$.  The triple $(V, \taufwd, \taubwd)$ is a
\emph{bitopological space} in the sense of Kelly~\cite{Kelly1963}.

The extension itself is a one-line corollary of the universality of generated topologies and carries essentially no mathematical content. The contribution of this paper begins only after this point. What is nontrivial is that this trivial extension canonically produces a bitopological structure whose pairwise invariants encode temporal asymmetry. The proof of this fact is structurally trivial: the topology axioms---closure
under finite intersections and arbitrary unions---are universal operations on
families of subsets of any set, and they are entirely indifferent to whether the
generating family arose from forward neighborhoods, backward neighborhoods, or
any other construction.  This structural triviality is exactly the point.  It
means that the extension from undirected to directed graphs is automatic and
requires no additional hypotheses (no symmetry, no connectedness, no acyclicity).
The mathematical content of the contribution lies not in the depth of the proof
but in identifying three things:
\begin{enumerate}[label=(\alph*),nosep]
  \item The asymmetry $\Nfwd{v} \neq \Nbwd{v}$ is precisely what captures
        temporal irreversibility: the topologies $\taufwd$ and $\taubwd$
        coincide if and only if $G$ is symmetric.
  \item The pair $(\taufwd, \taubwd)$ carries genuinely more information than
        either topology alone, because the pairwise separation properties of
        Kelly~\cite{Kelly1963} depend jointly on both topologies.
  \item The construction is tractable and yields quantitative invariants
        (component counts, base sizes, pairwise connectedness, clopen
        asymmetry) that can serve as features in time-series analysis.
\end{enumerate}

To verify the universality claim mechanically, we formalize the core
definitions and the central theorem in Lean~4 using the Mathlib library.
The formal proof reduces to two invocations of
\texttt{TopologicalSpace.generateFrom}, each of which satisfies the
topology axioms by construction, so the central theorem is mechanically
verified with no unproved placeholder.  A single \texttt{sorry} remains
elsewhere in the formalization, corresponding to the classical
McCord--Stong equivalence between topological connectedness and
comparability-graph connectedness for finite spaces
(\Cref{thm:specialization-components} below); this is a matter of
Mathlib infrastructure rather than mathematical uncertainty, since the
paper contains a complete classical proof of the theorem.  We discuss
the status of this gap in \Cref{subsec:limitations}.

The remainder of this paper is organized as follows.
\Cref{sec:prelim} establishes the necessary definitions and background
from point-set topology, graph theory, and the theory of bitopological
spaces.  \Cref{sec:extension} states and proves the central theorem,
establishes the relationship between the Nada topology and the
Alexandrov topology, defines irreversibility indices, and characterizes
pairwise connectedness.  \Cref{sec:algorithms} presents the algorithms
that realize the construction, together with polynomial-time complexity
analyses tied to the theorems of \Cref{sec:extension} and a three-valued
decision procedure for pairwise connectedness.  \Cref{sec:empirical}
applies the framework to the quarterly growth rate of U.S.\ real gross
domestic product over 1992--2024, obtaining explicit topological
invariants and a substantive finding about the COVID-19 shock.
\Cref{sec:discussion} discusses the relation to prior work, the
epistemic relation between the formalist proof and the constructive
algorithm, limitations, and directions for future work.
\Cref{sec:conclusion} summarizes.

% ══════════════════════════════════════════════════════════════════════════════
\section{Preliminaries}\label{sec:prelim}
% ══════════════════════════════════════════════════════════════════════════════

We establish the definitions and results that the central theorem depends on.
All sets are assumed to be subsets of a fixed nonempty set $V$ unless otherwise
stated.  We write $\powerset(V)$ for the powerset of $V$.

% ── 2.1 Topological spaces ──────────────────────────────────────────────────
\subsection{Topological spaces}\label{subsec:topology}

\begin{definition}[Topology]\label{def:topology}
A \emph{topology} on a set $X$ is a collection $\tau \subseteq \powerset(X)$
satisfying:
\begin{enumerate}[label=(\roman*),nosep]
  \item $\varnothing \in \tau$ and $X \in \tau$.
  \item If $U_1, U_2 \in \tau$, then $U_1 \cap U_2 \in \tau$
        (closure under finite intersections).
  \item If $\{U_i\}_{i \in I} \subseteq \tau$ for an arbitrary index set $I$,
        then $\bigcup_{i \in I} U_i \in \tau$ (closure under arbitrary unions).
\end{enumerate}
The pair $(X, \tau)$ is called a \emph{topological space}.  Elements of $\tau$
are called \emph{open sets}.  A set $C \subseteq X$ is \emph{closed} if
$X \setminus C \in \tau$.
\end{definition}

\begin{definition}[Subbase]\label{def:subbase}
A family $\subbase \subseteq \powerset(X)$ is a \emph{subbase} for a topology
$\tau$ on $X$ if the collection of all finite intersections of members of
$\subbase$ (including the empty intersection, which is $X$ by convention)
forms a \emph{base} for $\tau$.  Equivalently, $\tau$ is the coarsest topology
on $X$ containing every member of $\subbase$ as an open set.
\end{definition}

\begin{definition}[Base]\label{def:base}
A family $\base \subseteq \powerset(X)$ is a \emph{base} for a topology $\tau$
on $X$ if every open set $U \in \tau$ can be written as a union of members of
$\base$.  Equivalently, $\base$ is a base if and only if:
\begin{enumerate}[label=(\roman*),nosep]
  \item For every $x \in X$, there exists $B \in \base$ with $x \in B$.
  \item For every $B_1, B_2 \in \base$ and every $x \in B_1 \cap B_2$, there
        exists $B_3 \in \base$ with $x \in B_3 \subseteq B_1 \cap B_2$.
\end{enumerate}
\end{definition}

\begin{theorem}[Topology generated by an arbitrary family]\label{thm:generated-topology}
Let $X$ be any set and let $\subbase \subseteq \powerset(X)$ be any family of
subsets of $X$.  Then there exists a unique topology $\tau(\subbase)$ on $X$
that is the coarsest topology containing every member of $\subbase$.
Concretely, $\tau(\subbase)$ consists of all arbitrary unions of finite
intersections of members of $\subbase$, together with $\varnothing$ and $X$.

In particular, \emph{every} family of subsets of a set generates a topology.
No conditions on $\subbase$ are required.
\end{theorem}

\begin{proof}
The intersection of any collection of topologies on $X$ is again a topology on
$X$ (the axioms are preserved under arbitrary intersections of families of open
sets).  The discrete topology $\powerset(X)$ contains $\subbase$.  Therefore
the intersection of all topologies on $X$ that contain $\subbase$ is a
topology, and it is by construction the coarsest such topology.  The concrete
description follows by verifying that the family of arbitrary unions of finite
intersections of members of $\subbase$ satisfies the topology axioms.
\end{proof}

\begin{remark}\label{rem:universality}
\Cref{thm:generated-topology} is the key structural fact underlying the entire
paper.  The operations ``close under finite intersections'' and ``close under
arbitrary unions'' are operations on $\powerset(\powerset(X))$.  They are
\emph{indifferent} to how the family $\subbase$ was constructed---whether from
neighborhoods of an undirected graph, from forward neighborhoods of a digraph,
from open balls in a metric space, or from any other source.  This universality
is what makes the extension of the Nada procedure from undirected to directed
graphs automatic.
\end{remark}

% ── 2.2 Bitopological spaces ────────────────────────────────────────────────
\subsection{Bitopological spaces}\label{subsec:bitopology}

\begin{definition}[Bitopological space, Kelly~\cite{Kelly1963}]\label{def:bitopology}
A \emph{bitopological space} is a triple $(X, \tau_1, \tau_2)$ where $X$ is a
set and $\tau_1, \tau_2$ are both topologies on $X$.  No relationship between
$\tau_1$ and $\tau_2$ is assumed.
\end{definition}

\begin{definition}[Pairwise open and pairwise closed]\label{def:pairwise-open}
Let $(X, \tau_1, \tau_2)$ be a bitopological space.  A set $S \subseteq X$ is
\emph{pairwise open} if $S \in \tau_1 \cup \tau_2$, i.e., $S$ is open in at
least one of the two topologies.  A set $S$ is \emph{pairwise closed} if it is
closed in at least one of the two topologies.
\end{definition}

\begin{definition}[Pairwise Hausdorff]\label{def:pairwise-hausdorff}
A bitopological space $(X, \tau_1, \tau_2)$ is \emph{pairwise Hausdorff}
(or \emph{pairwise $T_2$}) if for every pair of distinct points $x, y \in X$,
there exist sets $U, V$ with $U \in \tau_1$, $V \in \tau_2$, $x \in U$,
$y \in V$, and $U \cap V = \varnothing$.
\end{definition}

\begin{definition}[Pairwise connected and pairwise disconnected]\label{def:pairwise-connected}
A bitopological space $(X, \tau_1, \tau_2)$ is \emph{pairwise disconnected} if
there exist nonempty sets $U \in \tau_1$ and $V \in \tau_2$ such that $U$ and
$V$ partition $X$ (i.e., $U \cup V = X$ and $U \cap V = \varnothing$).
It is \emph{pairwise connected} if it is not pairwise disconnected.
\end{definition}

\begin{remark}\label{rem:pairwise-vs-single}
The pairwise separation axioms of Kelly~\cite{Kelly1963} are genuinely stronger
than their single-topology counterparts.  A space may be connected with respect
to both $\tau_1$ and $\tau_2$ individually yet be pairwise disconnected, because
the partitioning sets may come from different topologies.  This is exactly the
mechanism by which the bitopological structure $(V, \taufwd, \taubwd)$ can
detect asymmetries invisible to either $\taufwd$ or $\taubwd$ alone.
\end{remark}

% ── 2.3 Graphs and neighborhoods ────────────────────────────────────────────
\subsection{Graphs and neighborhoods}\label{subsec:graphs}

\begin{definition}[Simple graph]\label{def:simple-graph}
A \emph{simple graph} is a pair $G = (V, E)$ where $V$ is a set (the
\emph{vertex set}) and $E \subseteq \binom{V}{2}$ is a set of unordered pairs
of distinct vertices (the \emph{edge set}).  The adjacency relation is
symmetric and irreflexive.
\end{definition}

\begin{definition}[Directed graph]\label{def:digraph}
A \emph{directed graph} (or \emph{digraph}) is a pair $G = (V, E)$ where $V$
is a set and $E \subseteq V \times V$ is a set of ordered pairs (the
\emph{arc set}).  The adjacency relation $\mathrm{Adj}(v, w) \iff (v, w) \in E$
is not required to be symmetric, reflexive, or irreflexive.
\end{definition}

\begin{definition}[Neighborhoods in undirected graphs]\label{def:undirected-nhd}
Let $G = (V, E)$ be a simple graph and $v \in V$.
\begin{itemize}[nosep]
  \item The \emph{(open) neighborhood} of $v$ is
        $N(v) = \{w \in V : \{v, w\} \in E\}$.
  \item The \emph{closed neighborhood} of $v$ is
        $\Ncl{v} = \{v\} \cup N(v)$.
\end{itemize}
\end{definition}

\begin{definition}[Forward and backward neighborhoods in digraphs]\label{def:directed-nhd}
Let $G = (V, E)$ be a digraph and $v \in V$.
\begin{itemize}[nosep]
  \item The \emph{forward (open) neighborhood} of $v$ is
        $N^{+}(v) = \{w \in V : (v, w) \in E\}$
        (the set of out-neighbors).
  \item The \emph{backward (open) neighborhood} of $v$ is
        $N^{-}(v) = \{w \in V : (w, v) \in E\}$
        (the set of in-neighbors).
  \item The \emph{forward closed neighborhood} of $v$ is
        $\Nfwd{v} = \{v\} \cup N^{+}(v)$.
  \item The \emph{backward closed neighborhood} of $v$ is
        $\Nbwd{v} = \{v\} \cup N^{-}(v)$.
\end{itemize}
\end{definition}

\begin{remark}\label{rem:symmetry-collapse}
If $G$ is symmetric (i.e., $(v,w) \in E \iff (w,v) \in E$ for all $v, w$),
then $\Nfwd{v} = \Nbwd{v}$ for every $v \in V$, and the forward and backward
constructions collapse to the undirected case.  The distinction between $\Nfwd{v}$
and $\Nbwd{v}$ is therefore a measure of the asymmetry of $G$.
\end{remark}

% ── 2.4 The Nada construction ────────────────────────────────────────────────
\subsection{The Nada construction for undirected graphs}\label{subsec:nada}

The following procedure was introduced by Nada et al.~\cite{Nada2018}.

\begin{definition}[Nada topology of an undirected graph]\label{def:nada-undirected}
Let $G = (V, E)$ be a simple graph.  The \emph{Nada topology} of $G$ is the
topology $\tau(G)$ on $V$ constructed as follows:
\begin{enumerate}[label=\textbf{Step \arabic*.},nosep]
  \item For each $v \in V$, compute the closed neighborhood
        $\Ncl{v} = \{v\} \cup N(v)$.
  \item Form the family $\subbase = \{\Ncl{v} : v \in V\}$.
  \item Close $\subbase$ under finite intersections to obtain a base
        $\base = \left\{\,\bigcap_{i=1}^{k} S_i : k \geq 0,\;
        S_1, \ldots, S_k \in \subbase\right\}$.
  \item Close $\base$ under arbitrary unions to obtain the topology
        \[
          \tau(G) \;=\; \Bigl\{\,\textstyle\bigcup_{j \in J} B_j \,:\,
          J \text{ arbitrary},\; B_j \in \base\Bigr\}.
        \]
\end{enumerate}
By \Cref{thm:generated-topology}, $\tau(G)$ is the unique coarsest topology
on $V$ containing every $\Ncl{v}$ as an open set.
\end{definition}

\begin{remark}\label{rem:nada-procedure-universal}
Steps 3 and 4 of \Cref{def:nada-undirected} are precisely the two operations in
\Cref{thm:generated-topology}.  The input to these steps is a family
$\subbase \subseteq \powerset(V)$, and the output is a topology
$\tau(\subbase)$.  Neither step inspects any property of $\subbase$ other than
its membership in $\powerset(V)$.  This is why the procedure extends to
directed graphs (or to any other source of subsets) without modification.
\end{remark}

% ── 2.5 Visibility graphs ───────────────────────────────────────────────────
\subsection{Visibility graphs}\label{subsec:visibility}

\begin{definition}[Natural Visibility Graph, Lacasa et al.~\cite{Lacasa2008}]%
\label{def:nvg}
Let $(x_1, x_2, \ldots, x_n)$ be a time series.  The \emph{Natural Visibility
Graph} (NVG) is the undirected graph $G = (V, E)$ where $V = \{1, 2, \ldots, n\}$
and two vertices $i < j$ are adjacent ($\{i, j\} \in E$) if and only if for
every intermediate index $i < k < j$,
\[
  x_k < x_j + (x_i - x_j)\,\frac{j - k}{j - i}\,.
\]
Geometrically, the data points $(i, x_i)$ and $(j, x_j)$ ``see'' each other if
the line segment connecting them passes above all intermediate data points.
\end{definition}

\begin{definition}[Horizontal Visibility Graph, Luque et al.~\cite{Luque2009}]%
\label{def:hvg}
The \emph{Horizontal Visibility Graph} (HVG) is defined as above, but with the
simpler condition: $\{i, j\} \in E$ if and only if $x_k < \min(x_i, x_j)$ for
every $i < k < j$.
\end{definition}

\begin{definition}[Directed Visibility Graphs]\label{def:directed-vg}
Given any visibility graph $G = (V, E)$ (NVG or HVG), the \emph{directed
visibility graph} is the digraph $\vec{G} = (V, \vec{E})$ obtained by orienting
every edge from the smaller index to the larger index: $(i, j) \in \vec{E}$ if
and only if $\{i, j\} \in E$ and $i < j$.

Since the index set $V = \{1, \ldots, n\}$ is totally ordered and edges are
oriented consistently with this order, $\vec{G}$ is a directed acyclic graph
(DAG) in which the time index serves as a topological order.
\end{definition}

\begin{remark}\label{rem:irreversibility}
In the directed visibility graph, the forward closed neighborhood $\Nfwd{i}$
consists of vertex $i$ and those later time points visible from~$i$, while the
backward closed neighborhood $\Nbwd{i}$ consists of vertex $i$ and those
earlier time points from which $i$ is visible.  For a time series generated by
a reversible process (e.g., an i.i.d.\ sequence), the statistical properties
of $\Nfwd{\cdot}$ and $\Nbwd{\cdot}$ are exchangeable.  For an irreversible
process (e.g., one with asymmetric relaxation), they are not.  The divergence
between $\taufwd$ and $\taubwd$ is therefore a topological signature of
temporal irreversibility.
\end{remark}

% ── 2.6 Alexandrov topology and specialization preorder ──────────────────────
\subsection{Alexandrov topology and the specialization preorder}%
\label{subsec:alexandrov}

For finite topological spaces, there is a classical correspondence between
topologies and preorders that provides efficient computational access to
topological invariants.

\begin{definition}[Alexandrov topology, Alexandrov~\cite{Alexandrov1937}]%
\label{def:alexandrov}
A topology $\tau$ on a set $X$ is an \emph{Alexandrov topology} if $\tau$ is
closed under arbitrary intersections (not just finite ones).  Equivalently,
every point $x \in X$ has a smallest open neighborhood.
\end{definition}

\begin{theorem}[Alexandrov correspondence for finite sets]%
\label{thm:alexandrov-correspondence}
Let $X$ be a finite set.  There is a bijection between topologies on $X$ and
preorders on $X$ (reflexive, transitive relations), given by:
\begin{itemize}[nosep]
  \item \textbf{Topology $\to$ preorder:}  Define $x \leq y$ if and only if
        $x \in \overline{\{y\}}$ (the closure of $\{y\}$), or equivalently,
        every open set containing $x$ also contains $y$.  This is the
        \emph{specialization preorder}.
  \item \textbf{Preorder $\to$ topology:}  The open sets are the \emph{upsets}
        (upper sets) of the preorder: a set $U \subseteq X$ is open if and only
        if $x \in U$ and $x \leq y$ implies $y \in U$.
\end{itemize}
Every topology on a finite set is an Alexandrov topology, so this
correspondence applies to all finite topological spaces.
\end{theorem}

\begin{proof}
Every topology on a finite set is Alexandrov because a finite intersection of
open sets is a finite intersection, which is open by the topology axioms, and
an arbitrary intersection of finitely many sets is a finite intersection.  The
bijection between Alexandrov topologies and preorders is classical;
see Alexandrov~\cite{Alexandrov1937}.
\end{proof}

\begin{definition}[Specialization preorder, McCord~\cite{McCord1966},
Stong~\cite{Stong1966}]\label{def:specialization}
Let $(X, \tau)$ be a topological space.  The \emph{specialization preorder}
$\leq_\tau$ is defined by
\[
  x \leq_\tau y \iff x \in \overline{\{y\}} \iff
  \text{every open set containing } x \text{ also contains } y.
\]
When $X$ is finite, $\leq_\tau$ is a preorder and the topology can be recovered
from it via the Alexandrov correspondence (\Cref{thm:alexandrov-correspondence}).
\end{definition}

\begin{theorem}[Connected components via the specialization preorder]%
\label{thm:connected-via-preorder}
Let $(X, \tau)$ be a finite topological space with specialization preorder
$\leq_\tau$.  Define the \emph{comparability graph} $C_\tau$ on $X$ by:
vertices $x$ and $y$ are adjacent in $C_\tau$ if and only if $x \leq_\tau y$
or $y \leq_\tau x$.  Then the connected components of $(X, \tau)$ (as a
topological space) coincide with the connected components of the graph $C_\tau$.
\end{theorem}

\begin{proof}
A self-contained proof is given in full for the logically equivalent
\Cref{thm:specialization-components} in
\Cref{sec:extension}; that proof establishes reflexivity and transitivity
of the specialization preorder, the identification of minimal open sets
in a finite Alexandrov space, and the coincidence between topological
connected components and comparability-graph connected components.  The
classical formulation of the statement is due to
McCord~\cite{McCord1966} and Stong~\cite{Stong1966}.
\end{proof}

\begin{remark}\label{rem:computational-significance}
\Cref{thm:connected-via-preorder} is computationally significant: it provides
an exact polynomial-time algorithm for computing the connected components of a
finite topological space without enumerating the topology.  One computes the
specialization preorder directly from the generating family (the subbase), then
finds connected components of the comparability graph using breadth-first
search.  This avoids the potentially exponential blowup of explicitly
constructing all open sets.  For the Nada construction, the specialization
preorder can be read off from the intersection structure of the closed
neighborhoods.
\end{remark}

% ══════════════════════════════════════════════════════════════════════════════
\section{Extension of Nada's Construction to Directed Graphs}%
\label{sec:extension}
% ══════════════════════════════════════════════════════════════════════════════

We now extend the Nada construction to directed graphs and prove that the
result is a bitopological space.  The extension requires no new mathematical
ideas: the universality of the generated-topology construction
(\Cref{thm:generated-topology}) does all the work.  The content of this section
is the identification of the directed neighborhoods as the appropriate input
to the Nada procedure, the analysis of the resulting bitopological structure,
and the relationship between the Nada topology and the Alexandrov topology on
the same vertex set.

% ── 3.1 Validity of the directed extension ──────────────────────────────────

\begin{theorem}[Validity of the directed extension]%
\label{thm:directed-validity}
Let $G = (V, E)$ be a directed graph.  The families
\[
  \subbase^{+} = \bigl\{\Nfwd{v} : v \in V\bigr\},
  \qquad
  \subbase^{-} = \bigl\{\Nbwd{v} : v \in V\bigr\}
\]
are valid subbases on~$V$, and the topologies $\taufwd$ and $\taubwd$ generated
from them by the Nada closure procedure (close the subbase under finite
intersections to obtain a base, then close the base under arbitrary unions to
obtain the topology) satisfy the topological space axioms.
\end{theorem}

\begin{proof}
By \Cref{thm:generated-topology}, every family $\subbase \subseteq \powerset(V)$
generates a unique topology $\tau(\subbase)$ on~$V$.  The closure operations
underlying the Nada procedure---finite intersection and arbitrary union in
$\powerset(V)$---are universal operations on families of subsets and do not
inspect any property of the generating family or the graph from which it arose.
Applying the theorem once to $\subbase^{+}$ yields $\taufwd$; applying it
once to $\subbase^{-}$ yields $\taubwd$.  Neither application requires
symmetry, connectedness, acyclicity, or any other hypothesis on~$G$.
\end{proof}

% ── 3.2 Self-inclusion guarantees non-degeneracy ────────────────────────────

\begin{proposition}[Self-inclusion guarantees non-degeneracy]%
\label{prop:self-inclusion}
The inclusion of $v$ in both $\Nfwd{v}$ and $\Nbwd{v}$---the ``self-bit''
$v \in \Nfwd{v}$ and $v \in \Nbwd{v}$---guarantees that no element of
$\subbase^{+}$ or $\subbase^{-}$ is empty.
\end{proposition}

\begin{proof}
For any vertex $v \in V$, the forward closed neighborhood is
$\Nfwd{v} = \{v\} \cup N^{+}(v)$, which contains $v$ regardless of whether
$N^{+}(v)$ is empty.  In particular, if $v$ is a sink (no out-edges), then
$N^{+}(v) = \varnothing$ and $\Nfwd{v} = \{v\} \neq \varnothing$.  Similarly,
if $v$ is a source (no in-edges), then $N^{-}(v) = \varnothing$ and
$\Nbwd{v} = \{v\} \neq \varnothing$.

Without the self-bit, the open out-neighborhood
$N^{+}(v) = \{w : (v,w) \in E\}$ of a sink vertex $v$ would be empty.  In a
non-trivial DAG (one with at least one vertex of out-degree zero, which every
finite DAG possesses), at least one element of the subbase $\{N^{+}(v) : v \in V\}$
would be $\varnothing$.  Since $\varnothing$ is contained in every finite
intersection, the intersection closure step would propagate $\varnothing$
throughout the base without adding information, collapsing the discriminative
power of the base: any difference between the neighborhood structure of
distinct sinks would be obliterated by the universal presence of the empty
set.  The self-bit prevents this degenerate collapse.
\end{proof}

% ── 3.3 Bitopological structure ─────────────────────────────────────────────

\begin{theorem}[Bitopological structure]%
\label{thm:bitopological-structure}
The triple $(V, \taufwd, \taubwd)$ is a bitopological space in the sense of
Kelly~\cite{Kelly1963}: both $\taufwd$ and $\taubwd$ are topologies on the
same underlying set~$V$.
\end{theorem}

\begin{proof}
By \Cref{thm:directed-validity}, $\taufwd$ and $\taubwd$ are both topologies
on~$V$.  A bitopological space is defined as a triple $(X, \tau_1, \tau_2)$
where $\tau_1$ and $\tau_2$ are topologies on $X$
(\Cref{def:bitopology}).  Setting $X = V$, $\tau_1 = \taufwd$,
$\tau_2 = \taubwd$ yields the result.
\end{proof}

% ── 3.4 Specialization preorder preserves connected components ──────────────

\begin{theorem}[Specialization preorder preserves connected components]%
\label{thm:specialization-components}
Let $(V, \tau)$ be a finite topological space with base~$\base$, and let
$\leq$ be the specialization preorder on $V$ defined by
$x \leq y$ if and only if $x \in U$ implies $y \in U$ for every $U \in \tau$
(equivalently, $\mathcal{S}(x) \subseteq \mathcal{S}(y)$, where
$\mathcal{S}(x) = \{B \in \base : x \in B\}$).  Define the undirected
comparability graph $C_\tau$ on $V$ by: vertices $x$ and $y$ are adjacent in
$C_\tau$ if and only if $x \leq y$ or $y \leq x$.  Then the topological
connected components of $(V, \tau)$ coincide exactly with the connected
components of $C_\tau$.  This applies to $\tau = \taufwd$ and
$\tau = \taubwd$ separately.
\end{theorem}

\begin{proof}
We establish the three parts in sequence.

\medskip
\noindent
\textbf{(a) The specialization preorder is a well-defined preorder.}
The specialization preorder $\leq_\tau$ on a topological space $(V, \tau)$ is
defined by $x \leq_\tau y$ if and only if every open set containing $x$ also
contains $y$ (\Cref{def:specialization}).  Reflexivity: every open set
containing $x$ contains $x$.  Transitivity: if every open set containing $x$
contains $y$, and every open set containing $y$ contains $z$, then every open
set containing $x$ contains $z$.  The equivalence classes of the symmetric
closure of $\leq_\tau$---i.e., the classes of the equivalence relation
$x \sim y \iff (x \leq y \text{ and } y \leq x)$---are exactly the fibers of
the quotient map $V \twoheadrightarrow V/{\sim}$, and $V/{\sim}$ equipped with
the quotient topology is the $T_0$ quotient (Kolmogorov quotient) of
$(V, \tau)$.  This is established by McCord~\cite{McCord1966} and
Stong~\cite{Stong1966}.

\medskip
\noindent
\textbf{(b) Topological connectedness coincides with comparability-graph connectedness.}
Every finite topological space is an Alexandrov space (\Cref{def:alexandrov}),
because a finite intersection of open sets is open by the topology axioms, and
an arbitrary intersection of finitely many sets is a finite intersection.
In an Alexandrov space, every point $x$ has a smallest open neighborhood
\[
  U(x) \;=\; \bigcap\{U \in \tau : x \in U\},
\]
which is open (as an arbitrary intersection of opens in an Alexandrov space).
The specialization preorder satisfies $x \leq y$ if and only if $y \in U(x)$
if and only if $U(x) \subseteq U(y)$.

\emph{Forward direction:} Suppose $x$ and $y$ lie in the same topological
connected component.  Since $(V, \tau)$ is finite, the connected component
$C$ is a connected subspace, and any open cover of $C$ by opens in $\tau$ has
a common point in the intersection of any two covering sets that overlap.  The
family $\{U(z) : z \in C\}$ covers $C$ (since $z \in U(z)$).  We construct
a chain $x = z_0, z_1, \ldots, z_k = y$ in $C$ such that consecutive points
are specialization-comparable, as follows.

Begin with $z_0 = x$.  The set $U(z_0)$ is a nonempty open subset of the
connected space $C$.  If $y \in U(z_0)$ or $z_0 \in U(y)$, we are done (the
chain has length~1).  Otherwise, since $C$ is connected and $\{U(z) : z \in C\}$
covers $C$, there is some $z_1 \in C$ with $z_1 \in U(z_0)$ or
$z_0 \in U(z_1)$, so $z_0 \leq z_1$ or $z_1 \leq z_0$.  Continue
inductively: the set of points reachable from $x$ by a chain of
specialization-comparable points is open (it is a union of minimal
neighborhoods) and closed (its complement is a union of minimal neighborhoods
of points not reachable from $x$).  Since $C$ is connected, this set is all
of $C$, so $y$ is reachable.

\emph{Backward direction:} Suppose $x = z_0, z_1, \ldots, z_k = y$ is a chain
of specialization-comparable points.  For each $i$, either $z_i \leq z_{i+1}$
or $z_{i+1} \leq z_i$.  In the first case, $z_{i+1} \in U(z_i)$, so both
$z_i$ and $z_{i+1}$ lie in the open connected set $U(z_i)$.  In the second
case, both lie in $U(z_{i+1})$.  In either case, $z_i$ and $z_{i+1}$ lie
in a common connected set, so they are in the same connected component.
By transitivity, $x$ and $y$ are in the same connected component.

\medskip
\noindent
\textbf{(c) Computational complexity.}
The set $\mathcal{S}(x) = \{B \in \base : x \in B\}$ can be represented as a
bitset over the base indices, requiring $O(|\base|/64)$ machine words per
vertex.  Constructing all $|V|$ bitsets requires $O(|V| \cdot |\base|)$
time and space.  Testing $x \leq y$ reduces to a bitset subset test:
$\mathcal{S}(x) \subseteq \mathcal{S}(y)$, which costs $O(|\base|/64)$ time
under word-parallel operations.  Constructing the comparability graph
$C_\tau$ requires $O(|V|^2)$ such tests, for a total of
$O(|V|^2 \cdot |\base|/64)$ time.  Finding connected components of $C_\tau$
by breadth-first search costs $O(|V|^2)$ time (since $C_\tau$ may have
$O(|V|^2)$ edges).  The total is $O(|V|^2 \cdot |\base|/64)$, which is
polynomial in $|V|$ and $|\base|$.  In particular, the connected components
of a finite topology can be decided without ever enumerating the full topology,
whose size may be exponential in $|V|$.
\end{proof}

% ── 3.5 Alexandrov topology as subtopology ──────────────────────────────────

\begin{theorem}[Alexandrov topology as subtopology of the directed Nada topology]%
\label{thm:alexandrov-subtopology}
Let $G = (V, E)$ be a finite DAG and let $R^*$ denote the
reflexive-transitive closure of $E$, a preorder on~$V$.  Let $\tau_A$ denote
the Alexandrov topology on $V$ whose open sets are exactly the upsets of
$R^*$.  Let $\taufwd_{\mathrm{Nada}}$ denote the forward Nada topology of
\Cref{thm:directed-validity}.  Then
\[
  \tau_A \;\subseteq\; \taufwd_{\mathrm{Nada}}.
\]
That is, every open set of the Alexandrov topology is also open in the
forward Nada topology.
\end{theorem}

\begin{proof}
Let $U$ be an upset of $R^*$.  We show $U$ is open in $\taufwd_{\mathrm{Nada}}$
by exhibiting it as a union of subbase elements.

For each $v \in U$, we claim that $\Nfwd{v} \subseteq U$.  To see this,
let $w \in \Nfwd{v}$.  Either $w = v$, in which case $w = v \in U$, or
$(v, w) \in E$, in which case $v \,R^*\, w$ (by a single-step path).  Since
$U$ is an upset of $R^*$ and $v \in U$, we conclude $w \in U$.

Since $v \in \Nfwd{v} \subseteq U$ for every $v \in U$, we have
\[
  U \;=\; \bigcup_{v \in U} \Nfwd{v}.
\]
The right-hand side is a union of elements of the subbase $\subbase^{+}$.
Every element of $\subbase^{+}$ is open in $\taufwd_{\mathrm{Nada}}$ (by
construction), and an arbitrary union of open sets is open.  Therefore $U$ is
open in $\taufwd_{\mathrm{Nada}}$.
\end{proof}

\begin{remark}[Base size versus topology inclusion]%
\label{rem:base-size-warning}
The inclusion $\tau_A \subseteq \taufwd_{\mathrm{Nada}}$ is a statement
about families of open sets, \emph{not} about the sizes of any
particular bases used to generate them.  It is a common but false
inference to conclude from $\tau_A \subseteq \taufwd_{\mathrm{Nada}}$
that the canonical base $\base_A$ of the Alexandrov topology has
cardinality at most that of the Nada base $\base^{+}_{\mathrm{Nada}}$.
A finer topology (one with more open sets) can in general be generated
by a \emph{smaller} base if the base is well chosen: base size is a
property of the generator, not of the topology.

In this paper, the empirical observation that the Nada base is typically
larger than the Alexandrov base on visibility graphs is a
computationally meaningful fact about the two specific pipelines used
to generate the topologies, not a corollary of
\Cref{thm:alexandrov-subtopology}.
\end{remark}

% ── 3.6 Resolution hierarchy ────────────────────────────────────────────────

\begin{corollary}[Resolution hierarchy]%
\label{cor:resolution-hierarchy}
If $\tau_A$ is connected and $\taufwd_{\mathrm{Nada}}$ is disconnected, then
the disconnection of $\taufwd_{\mathrm{Nada}}$ arises exclusively from open
sets in $\taufwd_{\mathrm{Nada}}$ that are not upsets of $R^*$, i.e., from
sets produced by the finite-intersection closure of the subbase $\subbase^{+}$
that do not belong to $\tau_A$.
\end{corollary}

\begin{proof}
Suppose for contradiction that every nontrivial clopen set of
$\taufwd_{\mathrm{Nada}}$ were also clopen in~$\tau_A$.  A clopen set of a
topology $\tau$ is a set $S$ with $S \in \tau$ and $S^c \in \tau$.  Since
$\taufwd_{\mathrm{Nada}}$ is disconnected, there exists a nonempty proper
subset $S \subsetneq V$ that is clopen in $\taufwd_{\mathrm{Nada}}$.  By
hypothesis, $S$ is also clopen in $\tau_A$.  But $S$ is nonempty and
$S \neq V$, so $S$ witnesses the disconnection of $\tau_A$: both $S$ and
$S^c$ are nonempty open sets in $\tau_A$ whose union is $V$ and whose
intersection is empty.  This contradicts the connectedness of $\tau_A$.

Therefore at least one disconnecting clopen of $\taufwd_{\mathrm{Nada}}$ is
not clopen in $\tau_A$.  Since $\tau_A \subseteq \taufwd_{\mathrm{Nada}}$
(\Cref{thm:alexandrov-subtopology}), any set that is open in $\tau_A$ is
automatically open in $\taufwd_{\mathrm{Nada}}$.  The additional open sets in
$\taufwd_{\mathrm{Nada}}$ that are not in $\tau_A$ arise precisely from the
finite-intersection closure of the subbase---the operation that takes
$\subbase^{+}$ and produces the base $\base^{+}$ by forming intersections
$\Nfwd{v_1} \cap \cdots \cap \Nfwd{v_k}$.  These intersection sets are the
source of the additional discriminative power of the Nada construction over
the pure Alexandrov construction on the same DAG.
\end{proof}

% ── 3.7 Irreversibility indices ─────────────────────────────────────────────

\begin{proposition}[Irreversibility indices]%
\label{prop:irreversibility-indices}
For a bitopological space $(V, \taufwd, \taubwd)$ arising from a directed
graph, define:
\begin{itemize}[nosep]
  \item the \emph{component-count irreversibility index}
        \[
          I_C(\taufwd, \taubwd)
          \;=\; \frac{|C^{+} - C^{-}|}{\max(C^{+}, C^{-})},
        \]
        where $C^{+}$ and $C^{-}$ are the numbers of connected components of
        $\taufwd$ and $\taubwd$ respectively, with the convention $I_C = 0$
        when $C^{+} = C^{-} = 0$;
  \item the \emph{base-size irreversibility index}
        \[
          I_B(\taufwd, \taubwd)
          \;=\; \frac{|\,|\base^{+}| - |\base^{-}|\,|}%
                     {\max(|\base^{+}|, |\base^{-}|)},
        \]
        with the same convention;
  \item the \emph{asymmetry direction}
        $\Delta(\taufwd, \taubwd) = C^{-} - C^{+}$,
        a signed integer.
\end{itemize}
Then:
\begin{enumerate}[label=(\roman*),nosep]
  \item $I_C, I_B \in [0, 1]$.
  \item $I_C = 0$ if and only if $C^{+} = C^{-}$, and $I_B = 0$ if and only
        if $|\base^{+}| = |\base^{-}|$.
  \item $\Delta > 0$ if and only if $\taufwd$ has strictly fewer connected
        components than $\taubwd$.
\end{enumerate}
\end{proposition}

\begin{proof}
(i) The numerator of $I_C$ is $|C^{+} - C^{-}| \leq \max(C^{+}, C^{-})$,
so $I_C \leq 1$.  Both numerator and denominator are nonneg\-ative, so
$I_C \geq 0$.  The same argument applies to $I_B$.

(ii) $I_C = 0$ iff $|C^{+} - C^{-}| = 0$ iff $C^{+} = C^{-}$, and
similarly for $I_B$.

(iii) $\Delta = C^{-} - C^{+} > 0$ iff $C^{+} < C^{-}$.
\end{proof}

\begin{remark}[Interpretation of the asymmetry direction]%
\label{rem:asymmetry-direction}
The sign of $\Delta$ indicates which of the two topologies is more
connected.  $\Delta > 0$ means $\taufwd$ has fewer connected components
than $\taubwd$, i.e., the forward topology is ``more connected'' than
the backward topology.  For directed visibility graphs of time series
with gradual rises and abrupt drops---the prototypical asymmetric
dynamics in the sense of Lacasa and Toral~\cite{LacasaToral2010}---the
prediction is $\Delta > 0$: forward visibility (looking ahead in time)
sees more structure than backward visibility (looking behind), because
the gradual rise creates many forward visibility edges while the abrupt
drop creates few backward edges.
\end{remark}

% ── 3.8 Characterization of pairwise connectedness ──────────────────────────

\begin{proposition}[Pairwise connectedness]%
\label{prop:pairwise-connected}
The bitopological space $(V, \taufwd, \taubwd)$ is pairwise connected in the
sense of Kelly~\cite{Kelly1963} if and only if there is no proper nonempty
subset $A \subsetneq V$ that is simultaneously $\taufwd$-open and
$\taubwd$-closed.  Equivalently, for every $U \in \taufwd$ with
$\varnothing \subsetneq U \subsetneq V$, the complement $V \setminus U$ does
not belong to $\taubwd$.
\end{proposition}

\begin{proof}
\emph{Equivalence of the two formulations.}  ``$A$ is $\taufwd$-open and
$\taubwd$-closed'' means $A \in \taufwd$ and $A^c \in \taubwd$.  The second
formulation states that for every $U \in \taufwd$ with
$\varnothing \subsetneq U \subsetneq V$, $U^c \notin \taubwd$.  Setting
$A = U$, the two conditions are identical.

\emph{Forward direction.}  Suppose $(V, \taufwd, \taubwd)$ is pairwise
connected.  By definition (\Cref{def:pairwise-connected}), there is no pair
of nonempty sets $U \in \taufwd$, $W \in \taubwd$ with $U \cup W = V$ and
$U \cap W = \varnothing$.  If some proper nonempty $A$ were $\taufwd$-open
and $\taubwd$-closed, then $U = A$ and $W = A^c$ would be such a pair (with
$W \in \taubwd$ since $A$ is $\taubwd$-closed), contradicting pairwise
connectedness.

\emph{Backward direction.}  Suppose no such $A$ exists, and suppose for
contradiction that $(V, \taufwd, \taubwd)$ is pairwise disconnected: there
exist nonempty $U \in \taufwd$ and $W \in \taubwd$ with $U \cup W = V$ and
$U \cap W = \varnothing$.  Then $W = U^c$, so $U$ is nonempty, proper (since
$W$ is nonempty), $\taufwd$-open, and $U^c = W \in \taubwd$ (i.e., $U$ is
$\taubwd$-closed).  This contradicts the hypothesis.

\medskip
\noindent
\emph{Computational remark.}  The characterization yields a decision procedure
that iterates over the enumerated open sets of $\taufwd$ and checks complement
membership in $\taubwd$, with total cost $O(|\taufwd| \cdot |\taubwd|)$ under
hash lookup.  When either topology has not been fully enumerated (because the
enumeration exceeded a resource bound \texttt{max\_open\_sets}), the decision
procedure returns a three-valued answer: pairwise connected, pairwise
disconnected, or \emph{undecided-within-budget}.  This three-valuedness is
epistemically honest rather than a defect of the algorithm: it reflects the
actual computational state after finite work.
\end{proof}

% ── Bibliography ─────────────────────────────────────────────────────────────
% ══════════════════════════════════════════════════════════════════════════════
\section{Algorithms and Computational Complexity}\label{sec:algorithms}
% ══════════════════════════════════════════════════════════════════════════════

Section~\ref{sec:extension} established the validity of the directed extension
and its structural properties.  We now describe the algorithms that realize
the construction and compute the bitopological invariants of
\Cref{prop:irreversibility-indices}, with explicit time and space complexity
bounds, termination arguments, and correctness justifications tied to the
theorems of Section~\ref{sec:extension}.  The algorithms are those implemented
in the reference R/C++ library accompanying this paper.  Each algorithm is
presented as formal pseudocode; each is accompanied by a proposition stating
its correctness and complexity, proved by direct reduction to a theorem of
Section~\ref{sec:extension}.

% ── 4.1 Bitset representation ───────────────────────────────────────────────

\subsection{Bitset representation of finite subsets}%
\label{subsec:bitset}

Let $V = \{0, 1, \ldots, n-1\}$.  Every subset $A \subseteq V$ is represented
as an array of $W = \lceil n/64 \rceil$ words of $64$ bits each; bit~$i$
(counting from the least significant bit of word $\lfloor i/64 \rfloor$) is
set if and only if $i \in A$.  The following operations are word-wise and
run in $O(W)$ time:
\begin{itemize}[leftmargin=2em]
  \item intersection $A \cap B$: word-wise bitwise \texttt{AND};
  \item union $A \cup B$: word-wise bitwise \texttt{OR};
  \item complement $V \setminus A$: word-wise bitwise \texttt{NOT} followed
        by masking the unused high bits of the last word;
  \item subset test $A \subseteq B$: verify that
        $A \cap \overline{B} = \varnothing$ word-wise;
  \item equality test $A = B$: word-wise comparison;
  \item emptiness test $A = \varnothing$: word-wise disjunction against zero;
  \item element test $i \in A$: single bit extraction.
\end{itemize}
All operations are branch-free in the inner loop and vectorize under modern
compilers.  A family of $k$ subsets stored consecutively occupies
$O(k \cdot W) = O(k \cdot n / 64)$ words in memory.

We also assume an auxiliary open-addressing hash table with word-wise hashing
(an \texttt{XOR}-fold of the $W$ words of each bitset), supporting amortized
$O(W)$-time insertion, membership query, and iteration.  The hash table is
used throughout to deduplicate families of bitsets as they are constructed.
We write $H_{\base}$, $H_{\tau}$, and so on for the hash tables associated
with particular families.

% ── 4.2 Closure under finite intersections ─────────────────────────────────

\subsection{Closure under finite intersections: wavefront propagation}%
\label{subsec:intersection-closure}

Given a subbase $\subbase = \{S_1, S_2, \ldots, S_k\} \subseteq \powerset(V)$,
the Nada construction requires computing the base
\[
  \base = \Bigl\{\bigcap_{i \in I} S_i : I \subseteq \{1, \ldots, k\},\,
  I \neq \varnothing,\, \bigcap_{i \in I} S_i \neq \varnothing\Bigr\},
\]
the closure of~$\subbase$ under nonempty finite intersections.  Naively
enumerating all $2^k - 1$ subsets of~$\subbase$ is infeasible.  We use a
\emph{wavefront} propagation: iteratively intersect a ``frontier'' of newly
discovered base elements with the existing base, retaining only new nonempty
intersections.

\begin{algorithm}
\caption{Intersection closure by wavefront propagation}%
\label{alg:intersection-closure}
\textbf{Input:} subbase $\subbase = \{S_1, \ldots, S_k\} \subseteq
\powerset(V)$, safety bound $M_{\base} \in \mathbb{N}$. \\
\textbf{Output:} the base~$\base$ together with a flag
\texttt{base\_complete} $\in \{\textup{true}, \textup{false}\}$.
\begin{enumerate}[leftmargin=2em, label=\textup{(\arabic*)}]
  \item Initialize $\base \leftarrow \subbase$ (deduplicated), insert every
        element of~$\subbase$ into a hash table $H_{\base}$, set the frontier
        $F \leftarrow \subbase$, and set \texttt{base\_complete}
        $\leftarrow$ \textup{true}.
  \item While $F \neq \varnothing$ and $|\base| \leq M_{\base}$:
  \begin{enumerate}[leftmargin=1.5em, label=\textup{(\alph*)}]
    \item Initialize $F' \leftarrow \varnothing$.
    \item For every $f \in F$ and every $b \in \base$ at the start of the
          iteration, compute $c \leftarrow f \cap b$.
    \item If $c \neq \varnothing$ and $c \notin H_{\base}$, insert $c$
          into~$H_{\base}$, append $c$ to~$\base$, and append $c$ to~$F'$;
          if $|\base| > M_{\base}$, set \texttt{base\_complete}
          $\leftarrow$ \textup{false} and break the loop.
  \end{enumerate}
  \item Set $F \leftarrow F'$ and continue.
  \item Return $\base$ and \texttt{base\_complete}.
\end{enumerate}
\end{algorithm}

\begin{proposition}[Correctness and complexity of
\Cref{alg:intersection-closure}]%
\label{prop:intersection-closure-correct}
\Cref{alg:intersection-closure} terminates in a finite number of iterations.
If \texttt{base\_complete} is \textup{true} at termination, then $\base$
equals the closure of~$\subbase$ under finite nonempty intersections.  In all
cases, $|\base| \leq M_{\base} + |\subbase|$ and every element of~$\base$ is
a finite nonempty intersection of elements of~$\subbase$.  The total time is
$O(|\base|^{2} \cdot W)$ word operations, and the space is
$O(|\base| \cdot W)$.
\end{proposition}

\begin{proof}
\emph{Termination.}  The family~$\base$ is a subset of $\powerset(V)$, which
has $2^n$ elements.  Each iteration of the main loop either strictly
increases~$|\base|$ or leaves it fixed; in the latter case $F' = \varnothing$
and the loop exits.  Hence the number of iterations is bounded by
$2^n - |\subbase|$.  A tighter practical bound: at iteration~$k$ (counting
from~$0$), every element of~$\base$ is an intersection of at most $k+1$
subbase elements, and any strictly descending chain of nonempty intersections
has length at most~$n$.  The number of iterations is therefore at most~$n$
once the frontier stabilizes.

\emph{Correctness.}  We show by induction on~$k$ that at the end of
iteration~$k$ of the main loop, $\base$ contains every nonempty intersection
$\bigcap_{i \in I} S_i$ with $|I| \leq k+1$.  The base case $k = 0$ is
step~(1): $\base$ is initialized to~$\subbase$, which contains every
single-element intersection.  For the inductive step, assume the claim holds
at the end of iteration~$k$.  At the start of iteration~$k+1$, the
frontier~$F$ consists precisely of the new intersections discovered in
iteration~$k$.  Let $T = \bigcap_{i \in I} S_i$ with $|I| = k+2$ be any
nonempty intersection of depth~$k+2$.  Choose any~$i_0 \in I$ and let
$T' = \bigcap_{i \in I \setminus \{i_0\}} S_i$.  Then $T'$ has depth~$k+1$,
so by the induction hypothesis $T' \in \base$ after iteration~$k$; and
$T = T' \cap S_{i_0}$ is an intersection of an element of~$\base$ with an
element of~$\subbase \subseteq \base$.  Step~(2b) of iteration~$k+1$
considers every pair $(f, b)$ with $f \in F \cup (\base \setminus F)$ and
$b \in \base$; in particular it considers $(T', S_{i_0})$ (in one of the two
orders), computes~$T$, and adds it to~$\base$ if not already present.
Therefore $\base$ contains every depth-$\leq k+2$ intersection at the end of
iteration~$k+1$.  Since every nonempty intersection of~$\subbase$ has depth
at most~$n$, the claim is established and~$\base$ stabilizes to the full
intersection closure at the latest after $n$ iterations.

\emph{Boundedness.}  Step~(2c) breaks the loop as soon as $|\base| >
M_{\base}$, so at termination $|\base| \leq M_{\base} + |\subbase|$ (the
additive $|\subbase|$ accounts for the initial population).  The
\texttt{base\_complete} flag records whether the bound was ever hit.

\emph{Complexity.}  Each iteration performs at most $|F| \cdot |\base|$
intersections, each costing $O(W)$ for the bitwise \texttt{AND} and $O(W)$
for the hash lookup.  Summing over iterations, the total number of
intersections is bounded by $\sum_{k} |F_k| \cdot |\base_k| \leq |\base|^2$,
yielding the stated bound of $O(|\base|^2 \cdot W)$.  The space is
$O(|\base| \cdot W)$ for the base and its hash table.
\end{proof}

\begin{remark}[Practical base sizes]\label{rem:practical-base-size}
For visibility graphs of time series, the empirical base sizes reported in
Section~\ref{sec:empirical} are polynomial in~$n$ and well within the
safety bound~$M_{\base}$.  The loose theoretical upper bound of~$2^n$ is
essentially never approached in practice because closed neighborhoods in
visibility graphs have substantial overlap, and the intersection lattice
collapses rapidly.
\end{remark}

% ── 4.3 Closure under arbitrary unions ─────────────────────────────────────

\subsection{Closure under arbitrary unions: wavefront propagation}%
\label{subsec:union-closure}

Given a base~$\base$, the Nada topology is
\[
  \tau = \Bigl\{\bigcup_{B \in \mathcal{F}} B : \mathcal{F} \subseteq \base
  \Bigr\} \cup \{\varnothing, V\},
\]
the closure of~$\base$ under arbitrary unions together with the empty set
and~$V$.  On a finite~$V$, arbitrary unions reduce to finite unions, and the
algorithm mirrors \Cref{alg:intersection-closure} exactly with the roles of
intersection and union exchanged.

\begin{algorithm}
\caption{Union closure by wavefront propagation}%
\label{alg:union-closure}
\textbf{Input:} base~$\base$, safety bound $M_{\tau} \in \mathbb{N}$. \\
\textbf{Output:} the topology~$\tau$ together with a flag
\texttt{topology\_complete}.
\begin{enumerate}[leftmargin=2em, label=\textup{(\arabic*)}]
  \item Initialize $\tau \leftarrow \base \cup \{\varnothing, V\}$, insert
        every element into a hash table~$H_{\tau}$, set the frontier
        $F \leftarrow \base$, and set \texttt{topology\_complete}
        $\leftarrow$ \textup{true}.
  \item While $F \neq \varnothing$ and $|\tau| \leq M_{\tau}$:
  \begin{enumerate}[leftmargin=1.5em, label=\textup{(\alph*)}]
    \item Initialize $F' \leftarrow \varnothing$.
    \item For every $f \in F$ and every $t \in \tau$ at the start of the
          iteration, compute $u \leftarrow f \cup t$.
    \item If $u \notin H_{\tau}$, insert~$u$ into~$H_{\tau}$, append to
          $\tau$ and to~$F'$; if $|\tau| > M_{\tau}$, set
          \texttt{topology\_complete} $\leftarrow$ \textup{false} and break.
  \end{enumerate}
  \item Set $F \leftarrow F'$ and continue.
  \item Return $\tau$ and \texttt{topology\_complete}.
\end{enumerate}
\end{algorithm}

\begin{proposition}[Correctness and complexity of \Cref{alg:union-closure}]%
\label{prop:union-closure-correct}
\Cref{alg:union-closure} terminates in a finite number of iterations.
If \texttt{topology\_complete} is \textup{true} at termination, then
$\tau$ equals the Nada topology generated by~$\base$.  In all cases,
$|\tau| \leq M_{\tau} + |\base| + 2$ and every element of~$\tau$ is
either $\varnothing$, $V$, or a finite union of elements of~$\base$.
The total time is $O(|\tau|^{2} \cdot W)$ word operations, and the
space is $O(|\tau| \cdot W)$.
\end{proposition}

\begin{proof}
\emph{Termination.}  The family~$\tau$ is a subset of $\powerset(V)$,
which has $2^{n}$ elements.  Each iteration of the main loop either
strictly increases $|\tau|$ or leaves it fixed; in the latter case
$F' = \varnothing$ and the loop exits.  The number of iterations is
therefore bounded by $2^{n}$, and a tighter bound follows from the
observation that at the end of iteration~$k$, every element of~$\tau$
is a union of at most $k + 1$ base elements; since any strictly
ascending chain of such unions in $\powerset(V)$ has length at
most~$n$, the number of iterations is at most~$n$ once the frontier
stabilizes.

\emph{Correctness.}  We show by induction on~$k$ that at the end of
iteration~$k$ of the main loop, $\tau$ contains every union
$\bigcup_{j \in J} B_j$ with $J \subseteq \{1, \ldots, |\base|\}$ and
$|J| \leq k + 1$.  The base case $k = 0$ is step~(1): $\tau$ is
initialized to $\base \cup \{\varnothing, V\}$, which contains every
single-element union (each $B_j$ by itself).  For the inductive step,
assume the claim holds at the end of iteration~$k$, and let
$T = \bigcup_{j \in J} B_j$ with $|J| = k + 2$.  Choose any
$j_{0} \in J$ and let $T' = \bigcup_{j \in J \setminus \{j_{0}\}} B_j$;
then $T'$ has depth $k + 1$, so by the induction hypothesis
$T' \in \tau$ at the start of iteration~$k + 1$.  Since $T'$ was
freshly added in iteration~$k$, it lies in the frontier~$F$ at the
start of iteration~$k + 1$.  Step~(2b) then considers the pair
$(T', B_{j_{0}})$, computes $T = T' \cup B_{j_{0}}$, and adds~$T$
to~$\tau$ if not already present.  Hence $\tau$ contains every
depth-$\leq\! k + 2$ union at the end of iteration~$k + 1$.  Because
every element of the Nada topology is a finite union of base elements
(plus $\varnothing$ and~$V$, already inserted in step~(1)), and every
such union has depth at most $|\base|$, the claim establishes that
$\tau$ stabilizes to the full Nada topology within at most $n$
iterations.

\emph{Boundedness and complexity.}  Step~(2c) breaks the loop as soon
as $|\tau| > M_{\tau}$, so at termination $|\tau| \leq M_{\tau} +
|\base| + 2$ (the additive term accounts for the initial population
with $\base \cup \{\varnothing, V\}$).  The
\texttt{topology\_complete} flag records whether the bound was hit.
Each iteration performs at most $|F| \cdot |\tau|$ unions, each
costing $O(W)$ for the bitwise \texttt{OR} and $O(W)$ for the hash
lookup.  Summing over iterations, the total number of union
operations is bounded by $\sum_{k} |F_{k}| \cdot |\tau_{k}| \leq
|\tau|^{2}$, yielding the stated bound of $O(|\tau|^{2} \cdot W)$
word operations.  The space is $O(|\tau| \cdot W)$ for $\tau$ and its
hash table.
\end{proof}

\begin{remark}
In principle $|\tau|$ can be as large as $2^{|\base|}$, making topology
enumeration by far the most expensive step of the pipeline.  The reference
implementation therefore uses a conservative default safety bound
$M_{\tau}$ and, for most workflows, enumerates the topology only when
pairwise connectedness (\Cref{prop:pairwise-connected}) is explicitly
requested.  The connected components of each marginal topology
$\taufwd$ and $\taubwd$, as well as the base-size asymmetry $\Delta$ and the
irreversibility indices $I_C, I_B$ of \Cref{prop:irreversibility-indices},
are computable directly from the \emph{base} without enumerating the
topology; see \Cref{subsec:specialization-algorithm} below.
\end{remark}

% ── 4.4 Specialization preorder algorithm ──────────────────────────────────

\subsection{Exact connected components via the specialization preorder}%
\label{subsec:specialization-algorithm}

\Cref{thm:specialization-components} reduces the computation of the
topologically connected components of a finite space $(V, \tau)$ to the
computation of the connected components of the comparability graph of the
specialization preorder, which in turn reduces to subset-inclusion tests of
the ``signature'' sets $\mathcal{S}(x) = \{B \in \base : x \in B\}$.  The
following algorithm realizes this reduction using the bitset representation,
operating on the base~$\base$ alone and therefore avoiding any dependence
on~$|\tau|$.

\begin{algorithm}
\caption{Connected components via the specialization preorder}%
\label{alg:specialization-components}
\textbf{Input:} base $\base = \{B_0, B_1, \ldots, B_{b-1}\}$ of a finite
topological space on~$V = \{0, \ldots, n-1\}$. \\
\textbf{Output:} the partition $V = C_1 \sqcup \cdots \sqcup C_r$ into
topologically connected components.
\begin{enumerate}[leftmargin=2em, label=\textup{(\arabic*)}]
  \item \emph{Signature construction.}  For each $x \in V$, allocate a bitset
        $\mathcal{S}(x)$ of length~$b$ (stored as $W_{\base} = \lceil b/64
        \rceil$ words initialized to zero).  For every $j \in \{0, \ldots,
        b-1\}$ and every $x \in B_j$, set bit~$j$ of~$\mathcal{S}(x)$.
  \item \emph{Comparability-graph traversal.}  Initialize $\mathrm{comp}(x)
        \leftarrow -1$ for all $x \in V$, and a component counter
        $r \leftarrow 0$.  For each $x \in V$ in order, if $\mathrm{comp}(x)
        = -1$, start a breadth-first search from~$x$ over the comparability
        graph, assigning $\mathrm{comp}(\cdot) \leftarrow r$ to every vertex
        reached, and then increment~$r$.  Two vertices $x$ and $y$ are
        adjacent in the comparability graph if and only if
        $\mathcal{S}(x) \subseteq \mathcal{S}(y)$ or
        $\mathcal{S}(y) \subseteq \mathcal{S}(x)$, tested word-wise on the
        bitset representations.
  \item Return the partition $\{C_1, \ldots, C_r\}$ where
        $C_i = \{x \in V : \mathrm{comp}(x) = i - 1\}$.
\end{enumerate}
\end{algorithm}

\begin{proposition}[Correctness and complexity of
\Cref{alg:specialization-components}]%
\label{prop:specialization-complexity}
\Cref{alg:specialization-components} returns the topologically connected
components of $(V, \tau)$ exactly, not as a heuristic approximation.  Its
time complexity is $O(n \cdot b + n^{2} \cdot W_{\base})$ and its space
complexity is $O(n \cdot W_{\base})$, where $b = |\base|$ and
$W_{\base} = \lceil b / 64 \rceil$.  In particular, the algorithm runs in
polynomial time in~$n$ and~$b$ without enumerating the topology, which may
be exponentially larger than~$b$.
\end{proposition}

\begin{proof}
Correctness is an immediate consequence of
\Cref{thm:specialization-components}: the topological connected components
of a finite space coincide with the connected components of the
comparability graph of the specialization preorder, and two vertices are
comparable in that preorder if and only if one signature is a subset of the
other.  The breadth-first search of step~(2) computes the connected
components of an undirected graph by definition.

For complexity, step~(1) performs
$\sum_{j=0}^{b-1} |B_j| = O(n \cdot b)$ bit assignments in the worst case
(each element of each base set contributes one bit).  Step~(2) performs at
most $O(n^{2})$ pairwise subset tests, each costing $O(W_{\base}) =
O(b / 64)$ word operations, for a total of $O(n^{2} \cdot b / 64)$.  Adding
the two contributions yields the stated bound.  The signature array
requires $O(n \cdot W_{\base})$ words of memory.
\end{proof}

By \Cref{prop:specialization-complexity}, connectedness---and therefore
the component-count irreversibility index~$I_C$, the asymmetry
direction~$\Delta$, and the base-size irreversibility index~$I_B$ of
\Cref{prop:irreversibility-indices}---is computable in polynomial time
from the base alone, without any dependence on~$|\tau|$.  In the reference implementation, the topology is enumerated
only when pairwise connectedness (\Cref{subsec:pairwise-algorithm}) is
explicitly requested.

% ── 4.5 Alexandrov topology via reverse-order propagation ──────────────────

\subsection{The Alexandrov topology via reverse-order reachability
propagation}%
\label{subsec:alexandrov-algorithm}

The directed Nada construction is complemented by the Alexandrov topology
induced by the reachability preorder of the DAG
(\Cref{thm:alexandrov-subtopology}).  For a DAG whose vertices are given in
topological order, this topology admits a particularly efficient computation
that exploits the order and requires only a single sweep.

Let $G = (V, E)$ be a finite DAG with $V = \{0, 1, \ldots, n-1\}$, satisfying
the topological-order property that every edge $(i, j) \in E$ has $i < j$.
(A directed visibility graph of a time series always satisfies this, with
the time index playing the role of topological order.)  Let
$R \subseteq V \times V$ denote the edge relation of~$G$ and let~$R^{*}$ be
its reflexive-transitive closure.  The following lemma is the key ingredient
that justifies computing the Alexandrov topology from single-step
reachability.

\begin{lemma}[Upsets of~$R$ coincide with upsets of~$R^{*}$]%
\label{lem:upsets-R-Rstar}
Let $R$ be any binary relation on a finite set~$V$, and let~$R^{*}$ be its
reflexive-transitive closure.  A subset $U \subseteq V$ is an upset of~$R$
(i.e., $x \in U$ and $x R y$ imply $y \in U$) if and only if~$U$ is an upset
of~$R^{*}$ (i.e., $x \in U$ and $x R^{*} y$ imply $y \in U$).
\end{lemma}

\begin{proof}
If $U$ is an upset of~$R^{*}$, then since $R \subseteq R^{*}$, any
$R$-successor of an element of~$U$ is also an $R^{*}$-successor, hence lies
in~$U$; so $U$ is an upset of~$R$.

Conversely, suppose~$U$ is an upset of~$R$, let $x \in U$, and let
$x R^{*} y$.  By definition of the reflexive-transitive closure, there
exists a finite chain
\[
  x = z_0 \; R \; z_1 \; R \; \cdots \; R \; z_m = y
\]
of length $m \geq 0$ (with $m = 0$ meaning $y = x \in U$ by reflexivity of
$R^{*}$).  Induction on~$m$ shows that every~$z_i$ lies
in~$U$: the base case $z_0 = x \in U$ is the hypothesis, and if
$z_i \in U$ then $z_i R z_{i+1}$ together with the upset property of~$U$
forces $z_{i+1} \in U$.  Therefore $y = z_m \in U$, proving that~$U$ is an
upset of~$R^{*}$.
\end{proof}

\Cref{lem:upsets-R-Rstar} tells us that the Alexandrov topology~$\tau_A$
with open sets equal to the upsets of~$R^{*}$ coincides with the topology
whose open sets are the upsets of~$R$.  We may therefore compute the
minimal open set of each vertex directly from the single-step relation,
provided we iterate along~$R$ once per vertex to accumulate multi-step
reachability.  The reverse topological order of the DAG allows this to be
done in a single sweep with bitset \texttt{OR}.

\begin{algorithm}
\caption{Alexandrov topology by reverse-order bitset propagation}%
\label{alg:alexandrov}
\textbf{Input:} a DAG $G = (V, E)$ with $V = \{0, \ldots, n-1\}$ and every
edge $(i, j) \in E$ satisfying $i < j$; the forward out-neighborhoods
$\mathrm{out}(v) = \{w : (v, w) \in E\}$. \\
\textbf{Output:} an array $\mathrm{reach}[\,\cdot\,]$ of bitsets over~$V$
such that $\mathrm{reach}[v]$ equals the set of~$R^{*}$-successors of~$v$,
which is the minimal open set of~$v$ in the Alexandrov topology~$\tau_A$.
\begin{enumerate}[leftmargin=2em, label=\textup{(\arabic*)}]
  \item For every $v \in V$, initialize the bitset $\mathrm{reach}[v]
        \leftarrow \{v\}$.
  \item For $v = n-1, n-2, \ldots, 0$ in reverse order: for every
        $w \in \mathrm{out}(v)$, set $\mathrm{reach}[v] \leftarrow
        \mathrm{reach}[v] \cup \mathrm{reach}[w]$ (one bitset \texttt{OR}).
  \item Return $\mathrm{reach}[\,\cdot\,]$.  The Alexandrov base is obtained
        by deduplicating $\{\mathrm{reach}[v] : v \in V\}$ via a hash table;
        the full topology~$\tau_A$, if required, is obtained by closing the
        base under arbitrary unions via \Cref{alg:union-closure}.
\end{enumerate}
\end{algorithm}

\begin{proposition}[Correctness and complexity of \Cref{alg:alexandrov}]%
\label{prop:alexandrov-correct}
After \Cref{alg:alexandrov} terminates, $\mathrm{reach}[v]$ equals the set
of all~$R^{*}$-successors of~$v$, which by \Cref{lem:upsets-R-Rstar} is the
minimal open set of~$v$ in~$\tau_A$.  The algorithm runs in time
$O(n + m \cdot n / 64)$ and uses space $O(n^{2} / 64)$, where
$m = |E|$.
\end{proposition}

\begin{proof}
\emph{Correctness.}  We prove by reverse induction on~$v$ that, after the
iteration of the main loop indexed by~$v$, $\mathrm{reach}[v]$ equals the
set of all~$R^{*}$-successors of~$v$.

\emph{Base case} ($v = n - 1$).  The topological-order property forces
$\mathrm{out}(n-1) = \varnothing$: any edge $(n-1, w) \in E$ would require
$w > n-1$, which is impossible because $n-1$ is the maximum index.  Hence
the inner loop is vacuous and $\mathrm{reach}[n-1] = \{n-1\}$, which equals
the set of $R^{*}$-successors of~$n-1$ because $n-1$ is a maximal element.

\emph{Inductive step.}  Suppose the claim holds for every $v' > v$.  At the
start of the iteration indexed by~$v$, we have $\mathrm{reach}[v] = \{v\}$
from the initialization in step~(1).  For every $w \in \mathrm{out}(v)$, the
topological-order property yields $w > v$, so by the induction hypothesis
$\mathrm{reach}[w]$ already equals the set of $R^{*}$-successors of~$w$.
The inner loop then sets
\[
  \mathrm{reach}[v] \;=\; \{v\} \;\cup\; \bigcup_{w \in \mathrm{out}(v)}
  \mathrm{reach}[w] \;=\; \{v\} \;\cup\; \bigcup_{w \in \mathrm{out}(v)}
  \bigl\{u : w R^{*} u\bigr\}.
\]
By definition of the reflexive-transitive closure, the set of
$R^{*}$-successors of~$v$ is exactly $\{v\}$ together with the union of the
sets of $R^{*}$-successors of all one-step $R$-successors of~$v$, which is
the right-hand side of the above display.  Hence $\mathrm{reach}[v]$ is
correct at the end of iteration~$v$.  Reverse induction on $v$ from~$n-1$
down to~$0$ establishes the claim for all vertices.

That the resulting family $\{\mathrm{reach}[v] : v \in V\}$ is a base for
the Alexandrov topology~$\tau_A$ follows from the standard correspondence
between preorders and Alexandrov topologies on finite sets
(\Cref{thm:alexandrov-correspondence}) together with
\Cref{lem:upsets-R-Rstar}: the minimal open set of~$v$ in any Alexandrov
topology is the intersection of all open sets containing~$v$, and for the
topology of upsets of~$R^{*}$ this minimal open set is precisely the
$R^{*}$-upward closure of~$\{v\}$, which equals $\mathrm{reach}[v]$.

\emph{Complexity.}  Step~(1) costs $O(n)$.  In step~(2), each outgoing edge
$(v, w) \in E$ contributes one bitset \texttt{OR} of cost $O(W) = O(n/64)$,
and each vertex is visited exactly once.  The total cost of the main loop
is $\sum_{v \in V} |\mathrm{out}(v)| \cdot O(n/64) = O(m \cdot n / 64)$.
Adding the initialization cost yields the stated time bound.  The space is
dominated by the $\mathrm{reach}$ array, which holds~$n$ bitsets each of
$O(n/64)$ words, for a total of $O(n^{2}/64)$.
\end{proof}

\begin{remark}[Optimality]\label{rem:alexandrov-optimal}
\Cref{alg:alexandrov} is asymptotically optimal for dense DAGs under the
word-RAM model with $64$-bit words: no algorithm can compute the transitive
closure in asymptotically less than $\Theta(n \cdot m / 64)$ time in this
model.  The reverse-order sweep achieves this bound while simultaneously
exploiting the topological order to avoid any explicit stabilization loop.
The Alexandrov base produced by \Cref{alg:alexandrov} can be fed directly
into \Cref{alg:specialization-components} to obtain the topological
connected components of~$\tau_A$ in the polynomial time bound of
\Cref{prop:specialization-complexity}, once again without enumerating
$\tau_A$ itself.
\end{remark}

% ── 4.6 Pairwise connectedness ─────────────────────────────────────────────

\subsection{Pairwise connectedness and the three-valued return}%
\label{subsec:pairwise-algorithm}

The pairwise connectedness test of \Cref{prop:pairwise-connected} is the
only operation in the pipeline that requires the full enumeration of both
topologies $\taufwd$ and $\taubwd$.  It is therefore the most expensive
step and the one most likely to exceed the safety bound~$M_{\tau}$ of
\Cref{alg:union-closure}.  This section describes the algorithm and gives
particular attention to the constructive handling of budget overflow.

\begin{algorithm}
\caption{Pairwise connectedness check}%
\label{alg:pairwise}
\textbf{Input:} enumerated topologies $\taufwd$ and $\taubwd$ on~$V$, each
with its \texttt{topology\_complete} flag from \Cref{alg:union-closure}. \\
\textbf{Output:} one of the three values $\{\textup{pairwise connected},
\textup{pairwise disconnected}, \textup{undecided}\}$, together with a
witness pair $(U, V \setminus U)$ when the answer is \textup{pairwise
disconnected}.
\begin{enumerate}[leftmargin=2em, label=\textup{(\arabic*)}]
  \item Build a hash table $H^{-}$ containing every $W \in \taubwd$ with
        $\varnothing \neq W \neq V$.
  \item For every $U \in \taufwd$ with $\varnothing \neq U \neq V$, compute
        $W \leftarrow V \setminus U$ (bitset complement with masking of the
        unused high bits of the last word).  If $W \in H^{-}$, return
        \textup{pairwise disconnected} with the witness pair $(U, W)$.
  \item If the loop of step~(2) completes without finding a witness, then:
  \begin{itemize}[leftmargin=1.5em]
    \item if both \texttt{topology\_complete} flags are \textup{true},
          return \textup{pairwise connected};
    \item otherwise, return \textup{undecided}.
  \end{itemize}
\end{enumerate}
\end{algorithm}

\begin{proposition}[Soundness, complexity, and epistemic status of
\Cref{alg:pairwise}]%
\label{prop:pairwise-correct}
\Cref{alg:pairwise} is sound with respect to the pairwise connectedness
predicate of \Cref{def:pairwise-connected}.  Specifically:
\begin{enumerate}[leftmargin=2em, label=\textup{(\roman*)}]
  \item if it returns \textup{pairwise disconnected} with witness $(U, W)$,
        then $U \in \taufwd$, $W \in \taubwd$, $U \cup W = V$, and
        $U \cap W = \varnothing$, so $(U, W)$ genuinely witnesses the
        disconnection;
  \item if it returns \textup{pairwise connected}, the bitopological space
        $(V, \taufwd, \taubwd)$ is pairwise connected;
  \item if it returns \textup{undecided}, nothing is claimed about the
        pairwise connectedness of $(V, \taufwd, \taubwd)$.
\end{enumerate}
The time complexity is $O(|\taufwd| \cdot W + |\taubwd| \cdot W)$ and the
space complexity is $O(|\taubwd| \cdot W)$ for the hash table.
\end{proposition}

\begin{proof}
Claim~(i): the hash table~$H^{-}$ is built from~$\taubwd$, so any~$W$ that
it contains lies in~$\taubwd$.  The set $W = V \setminus U$ by
construction, so $U \cup W = V$ and $U \cap W = \varnothing$ identically.
Hence the witness pair is genuine.

Claim~(ii): a \textup{pairwise connected} return requires both
\texttt{topology\_complete} flags to be \textup{true}, so the loop of
step~(2) has examined every proper nonempty $U \in \taufwd$ and every
proper nonempty $W \in \taubwd$ is present in~$H^{-}$.  The failure to
find a witness means no such pair $(U, W)$ with $W = V \setminus U$ exists
in $\taufwd \times \taubwd$, which by
\Cref{prop:pairwise-connected} is the definition of pairwise
connectedness.

Claim~(iii) is immediate: the \textup{undecided} branch makes no claim.
Complexity is dominated by the single pass over $\taufwd$ in step~(2),
each iteration performing one bitset complement and one hash lookup, each
$O(W)$.  The hash table occupies $O(|\taubwd| \cdot W)$ space.
\end{proof}

The three-valued return of \Cref{alg:pairwise} deserves explicit discussion
because it encodes a methodological commitment.  When either
\texttt{topology\_complete} flag is \textup{false}, \Cref{alg:pairwise}
has not completed the exhaustive search over the full~$\taufwd$, or has
compared against an incomplete~$H^{-}$ (or both), and therefore cannot
soundly claim pairwise connectedness.  A classical implementation would be
tempted to treat budget overflow as ``good enough'' and default to one of
the two Boolean answers; but either choice would attach a truth value to a
proposition that has not been decided with the resources consumed.

The constructively correct response is to return a third value,
\textup{undecided}, which reflects the actual epistemic state of the
computation after finite work.  The distinction between
``decided~\textup{true}'', ``decided~\textup{false}'', and ``not decided
within the available budget'' is not cosmetic: any downstream consumer of
the pairwise connectedness result must propagate \textup{undecided} as a
third value and cannot collapse it to a Boolean without losing soundness.
In practice, the \textup{undecided} return is a signal either to increase
$M_{\tau}$ and retry, or to rely on the polynomial-time marginal
invariants of \Cref{alg:specialization-components}---which do not require
topology enumeration---instead.  This design decision is the operational
counterpart, at the algorithmic level, of the constructivist discipline
articulated in Section~\ref{sec:intro}: a topological statement counts as
established only when the computation has finished the work that
establishes it.

% ── 4.7 Template dispatch ─────────────────────────────────────────────────

\subsection{Compile-time template dispatch for small ground sets}%
\label{subsec:template-dispatch}

The bitset operations of \Cref{subsec:bitset} all run in~$O(W)$, but the
constant hidden by the asymptotic notation depends strongly on whether the
loop over the $W$ words is known to be a compile-time constant or a runtime
variable.  When $W$ is known at compile time, modern C++ compilers fully
unroll the inner loop and the intersection, union, subset test, and
equality test reduce to a handful of machine instructions with zero loop
overhead.  When $W$ is a runtime variable, the inner loop incurs the cost
of loop bookkeeping, branch prediction, and potential pipeline stalls.
For the small ground sets typical of visibility graphs of short time
series, the compile-time case is overwhelmingly the common one, and it is
worth specializing.

The reference implementation uses C++ templates to dispatch at compile time
on three specialized values of~$W$, selected to cover the most common size
regimes for time-series applications:
\begin{itemize}[leftmargin=2em]
  \item $W = 1$, covering $n \leq 64$, in which every bitset operation
        becomes a single \texttt{uint64\_t} instruction;
  \item $W = 2$, covering $65 \leq n \leq 128$, in which operations are two
        instructions fully unrolled;
  \item $W = 3$, covering $129 \leq n \leq 192$, unrolled to three
        instructions.
\end{itemize}
For $n > 192$, the implementation falls back to a runtime loop over a
variable~$W$, which is slower per operation but imposes no upper bound
on~$n$.  The dispatch happens in a single switch statement at the entry
point of each algorithm, selecting the appropriate template specialization
based on $W = \lceil n / 64 \rceil$.

The template dispatch changes the constant factor but not the asymptotic
complexity of any algorithm.  All the bounds stated in
\Cref{prop:intersection-closure-correct},
\Cref{prop:union-closure-correct},
\Cref{prop:specialization-complexity},
\Cref{prop:alexandrov-correct}, and \Cref{prop:pairwise-correct} hold in
either dispatch regime; the specializations merely reduce the constant
hidden by $O(W)$ to a small fixed integer number of instructions for the
common small-$n$ case.  The practical speedup on the size regimes that
dominate time-series applications is substantial: between threefold and
tenfold over an unspecialized loop, depending on the operation and the
surrounding hot path.


% ══════════════════════════════════════════════════════════════════════════════
\section{Empirical Application: Topological Analysis of U.S.\ Real GDP Growth}%
\label{sec:empirical}
% ══════════════════════════════════════════════════════════════════════════════

We now apply the directed Nada bitopological construction of
\Cref{sec:extension}, together with the algorithms of
\Cref{sec:algorithms}, to a time series with well-known dynamical
asymmetries: the quarter-on-quarter growth rate of U.S.\ real gross
domestic product. The goal is twofold. First, to demonstrate that the
pipeline produces quantitative invariants on a real economic dataset,
with no tuning of free parameters and at a computational cost negligible
on commodity hardware. Second, and more importantly, to show that the
three topological layers built by the pipeline---the Alexandrov topology
$\tau_A$, the undirected Nada topology $\tau_{\text{Nada}}$, and the
directed Nada bitopology $(\taufwd, \taubwd)$---reveal \emph{distinct
structural facts} about the same series, each at a different level of
resolution, and that the asymmetry between $\taufwd$ and $\taubwd$
provides a purely topological quantification of temporal irreversibility
consistent with the classical economic literature on business-cycle
asymmetries.

%-------------------------------------------------------------------------------
\subsection{Data and construction pipeline}\label{subsec:data}
%-------------------------------------------------------------------------------

The series analyzed is the seasonally adjusted quarterly growth rate of
U.S.\ real gross domestic product, spanning Q1~1992 to Q1~2024 for a
total of $n = 129$ observations. The choice of this window is deliberate:
it straddles three distinct macroeconomic regime shifts---the
dot-com bust of 2001, the Global Financial Crisis of 2008--2009, and
the COVID-19 shock of 2020---while remaining within a single statistical
regime for the underlying measurement methodology. The combined presence
of these three episodes, each with markedly different speed and depth of
contraction and recovery, makes the series a canonical case of an
asymmetric dynamical system: expansions are typically gradual and
sustained, while contractions can be sharp or catastrophic. The
literature on business-cycle asymmetries, from Neftçi to Hamilton's
regime-switching models, has documented this temporal irreversibility
through statistical methods; here we approach it topologically.

The construction pipeline proceeds in the following stages:
\begin{enumerate}[label=(\roman*), leftmargin=2em]
\item From the raw series we construct the Horizontal Visibility Graph
(HVG) and the Natural Visibility Graph (NVG) in their directed form,
with edges oriented from earlier to later time indices. Both graphs are
DAGs on the same vertex set $V = \{1, 2, \ldots, 129\}$.
\item For each graph we compute the undirected Nada topology (closing
the family of undirected closed neighborhoods under finite intersection
and arbitrary union), the forward Nada topology $\taufwd$ (from the
forward closed neighborhoods $\Nfwd{v}$), the backward Nada topology
$\taubwd$ (from the backward closed neighborhoods $\Nbwd{v}$), and the
Alexandrov topology $\tau_A$ (from the reachability upsets of the
directed edge relation).
\item For each of these topologies we compute the exact number of
connected components via the specialization-preorder algorithm
(\Cref{alg:specialization-components}), together with the base size
$|\base|$ produced by the intersection-closure step
(\Cref{alg:intersection-closure}). For the Alexandrov topology we use
the reverse-order propagation of
\Cref{alg:alexandrov}.
\item From the forward and backward invariants we compute the three
bitopological quantities of Proposition~\ref{prop:irreversibility-indices}:
the component-count irreversibility index~$I_C$, the base-size
irreversibility index~$I_B$, and the asymmetry direction~$\Delta = C^-
- C^+$.
\end{enumerate}
The entire pipeline runs in well under one second on commodity hardware
for $n = 129$; the dominant cost is the intersection closure of the base,
which for the Horizontal Visibility Graph produces a base of size~$470$
and for the Natural Visibility Graph a base of size~$415$. No safety
bound is triggered for this~$n$; all topologies relevant to connectedness
are computed exactly, with complete bases.

%-------------------------------------------------------------------------------
\subsection{Results}\label{subsec:results}
%-------------------------------------------------------------------------------

\Cref{tab:gdp-results} collects the principal invariants for both graph
constructions and all four topologies. The entries are produced directly
by the reference implementation and not rounded beyond what is shown.

\begin{table}[h]
\centering
\small
\begin{tabular}{l c c c c c}
\hline
\textbf{Topology} & \textbf{Edges} & \textbf{Base size} &
  \textbf{Components} & \textbf{Connected?} & $\Delta / I_{C}$ \\
\hline
\multicolumn{6}{l}{\emph{Horizontal Visibility Graph (HVG)}} \\
Undirected Nada $\tau_{\text{Nada}}$ & $248$ & $470$ & $36$ & no & --- \\
Forward Nada $\taufwd$               & ---  & $250$ & $63$ & no & --- \\
Backward Nada $\taubwd$              & ---  & $250$ & $67$ & no & --- \\
Alexandrov $\tau_A$                  & ---  & $129$ & $\phantom{0}1$ & yes & --- \\[0.4ex]
\multicolumn{5}{l}{\hspace{1em}Asymmetry direction $\Delta = C^{-} - C^{+}$:}
  & $+4$ \\
\multicolumn{5}{l}{\hspace{1em}Component-count irreversibility $I_{C}$:}
  & $0.060$ \\
\multicolumn{5}{l}{\hspace{1em}Base-size irreversibility $I_{B}$:}
  & $0.000$ \\
\hline
\multicolumn{6}{l}{\emph{Natural Visibility Graph (NVG)}} \\
Undirected Nada $\tau_{\text{Nada}}$ & $406$ & $415$ & $\phantom{0}6$ & no & --- \\
Forward Nada $\taufwd$               & ---  & $239$ & $27$ & no & --- \\
Backward Nada $\taubwd$              & ---  & $239$ & $31$ & no & --- \\
Alexandrov $\tau_A$                  & ---  & $129$ & $\phantom{0}1$ & yes & --- \\[0.4ex]
\multicolumn{5}{l}{\hspace{1em}Asymmetry direction $\Delta = C^{-} - C^{+}$:}
  & $+4$ \\
\multicolumn{5}{l}{\hspace{1em}Component-count irreversibility $I_{C}$:}
  & $0.129$ \\
\multicolumn{5}{l}{\hspace{1em}Base-size irreversibility $I_{B}$:}
  & $0.000$ \\
\hline
\end{tabular}
\caption{Topological invariants of the directed Nada bitopological
analysis of U.S.\ real GDP quarterly growth, Q1~1992 to Q1~2024
($n = 129$). Edges refer to the undirected graph; forward and backward
topologies inherit the underlying directed edges by orientation.
Components are exact topological connected components computed via the
specialization preorder (\Cref{alg:specialization-components}). The
irreversibility indices $\Delta$ and $I_C$ are computed from the forward
and backward component counts as defined in
Proposition~\ref{prop:irreversibility-indices}; $I_B$ is computed from
the forward and backward base sizes.\label{tab:gdp-results}}
\end{table}

We now read \Cref{tab:gdp-results} line by line, distinguishing the
structural facts revealed by each topological layer.

\paragraph{The Alexandrov layer: global connectedness.}
For both graph constructions, the Alexandrov topology $\tau_A$ has
exactly one connected component and a base of size~$n = 129$. This is
the maximum possible connectedness of the topology, and it means that
the reachability preorder of the directed visibility graph links every
observation to every other observation through a chain of directed
paths. In economic terms, the U.S.\ GDP series is globally connected at
the level of \emph{causal-temporal precedence}: every quarter influences
every subsequent quarter through the compositional structure of the
business cycle. The base size of~$129$ equals the number of observations,
meaning that the minimal Alexandrov base consists of one minimal open
set per vertex and no coalescence occurs. This layer, on its own, would
lead to the conclusion that the system is indecomposable---a conclusion
that is true at the level of reachability but hides the finer structure
revealed by the higher layers.

\paragraph{The undirected Nada layer: local fragmentation.}
The picture changes drastically at the level of the undirected Nada
topology. For the HVG, $\tau_{\text{Nada}}$ has $36$ connected
components and a base of $470$ elements---more than three times the
number of vertices. For the NVG, it has $6$ connected components and a
base of $415$ elements. In both cases the topology is
\emph{disconnected}, which means that the closure of the closed
neighborhoods under finite intersection produces a partition of the
vertex set into topologically incommensurable clusters. The intersection
closure captures something that the pure reachability structure cannot:
local commonalities of neighborhood, where a single vertex may be linked
to two others that nevertheless belong to different Nada components
because their intersections with other neighborhoods isolate them. The
contrast with the Alexandrov layer is sharp: the same DAG is globally
connected under the reachability topology and highly fragmented under
the intersection-closed topology.

\paragraph{The bitopological layer: temporal asymmetry.}
The forward and backward Nada topologies yield strictly different
component counts in both graph constructions. For the HVG, $\taufwd$ has
$63$ components and $\taubwd$ has $67$; for the NVG, $\taufwd$ has $27$
components and $\taubwd$ has $31$. In both cases $\Delta = C^{-} - C^{+}
= +4$: the forward topology is strictly more connected than the backward
topology, by exactly four components. The component-count irreversibility
indices are $I_C = 0.060$ for the HVG and $I_C = 0.129$ for the NVG:
small in absolute value but strictly positive in both cases, and both
pointing in the same direction. Notably, the base-size irreversibility
index $I_B$ is \emph{exactly zero} in both constructions, since
$|\base^{+}| = |\base^{-}|$ in each graph ($250$ and $250$ for the HVG,
$239$ and $239$ for the NVG). The asymmetry between $\taufwd$ and
$\taubwd$ therefore lives \emph{entirely in the structure of the base
elements and their intersection closure}, not in the cardinality of the
generators. This is a nontrivial finding: the two subbases $\{\Nfwd{v}\}$
and $\{\Nbwd{v}\}$ produce the same number of base elements, but the
intersections of those elements distribute differently between the
forward and backward cases, yielding different component structures for
the specialization preorder.

The sign of $\Delta$ is consistent with the prediction of Lacasa and
Toral~\cite{LacasaToral2010} for time series with asymmetric dynamics:
in a regime of slow expansions and abrupt contractions, the forward
visibility is less obstructed than the backward visibility, so the
forward topology should be more connected (fewer components) than the
backward one. U.S.\ GDP in the 1992--2024 window is a paradigmatic
example of such dynamics: the expansions of 1992--2000, 2001--2007, and
2009--2020 are all gradual and monotone on aggregate, while the
contractions associated with the dot-com bust, the Global Financial
Crisis, and the COVID-19 shock are sharp and short-lived. Our
topological quantification of this asymmetry is the positive
$\Delta = +4$ in both HVG and NVG constructions. The numerical
invariance of $\Delta$ across the two graph constructions, despite the
markedly different edge counts and base sizes of HVG and NVG, suggests
that the asymmetry is a robust feature of the underlying series rather
than an artifact of the visibility-graph construction.

%-------------------------------------------------------------------------------
\subsection{Resolution gain: an empirical observation}%
\label{subsec:resolution-gain}
%-------------------------------------------------------------------------------

A striking feature of \Cref{tab:gdp-results} is the large gap between
the base size of the Nada topology and that of the Alexandrov topology.
For the HVG, $|\base_{\text{Nada}}| - |\base_A| = 470 - 129 = +341$,
which corresponds to a relative gain of
$341 / 470 \approx 72.6\%$. For the NVG, $|\base_{\text{Nada}}| -
|\base_A| = 415 - 129 = +286$, a relative gain of
$286 / 415 \approx 68.9\%$. In both cases, more than two thirds of the
base elements generated by the Nada intersection closure are \emph{not}
upsets of the reachability preorder.

\begin{remark}[Resolution gain is empirical, not a corollary]%
\label{rem:resolution-gain-empirical}
It is essential to read this observation correctly. The inclusion
$\tau_A \subseteq \taufwd_{\text{Nada}}$ established in
Theorem~\ref{thm:alexandrov-subtopology} is an inclusion of
\emph{topologies}, not of base sizes. A finer topology can in general be
generated by a \emph{smaller} base if the base is well chosen: base size
is a property of the generator, not of the topology it generates. The
observation that the Nada base is larger than the Alexandrov base on
the visibility graphs of this series is therefore a fact about the two
specific algorithms used to compute the generators---the intersection
closure of closed neighborhoods versus the deduplicated family of
reachability upsets---and \emph{not} a logical consequence of the
inclusion of topologies. We present the resolution gain as an empirical
quantity that measures the \emph{additional discriminative power}
extracted by the Nada intersection closure beyond the pure order
structure of the DAG, on the specific inputs analyzed here. On other
inputs, or with different generator conventions, the direction of the
gap could in principle be reversed, even though the inclusion of
topologies would continue to hold.
\end{remark}

With this caveat, the resolution gain on U.S.\ GDP is large and
consistent across the two graph constructions, which suggests that the
intersection closure extracts genuine topological information in this
case. Corollary~\ref{cor:resolution-hierarchy} then explains the
mechanism: because $\tau_A$ is connected and $\tau_{\text{Nada}}$ (both
undirected and directional) is disconnected, the disconnection of the
Nada topology arises entirely from open sets that are generated by
finite intersections of closed neighborhoods and that are \emph{not}
upsets of the reachability preorder. The Nada construction sees a
structure that the Alexandrov construction cannot see by design.

%-------------------------------------------------------------------------------
\subsection{A substantive finding: the COVID-19 shock}%
\label{subsec:covid-finding}
%-------------------------------------------------------------------------------

Among the $129$ observations, those at indices $114$ and $115$---%
corresponding to Q2 and Q3 of 2020, respectively---are the two most
extreme values of the entire series. The second quarter of 2020
records the sharpest contraction in U.S.\ real GDP on record, at an
annualized rate of approximately $-31\%$; the third quarter of the
same year records the sharpest rebound, at approximately $+33\%$
annualized. These two quarters are, by almost any descriptive
statistic, the outliers of the series.

The following observation, which emerges directly from the output of
\Cref{alg:specialization-components}, is worth making explicit: under
the Nada topology (both undirected and directed, both HVG and NVG),
observations~$114$ and~$115$ belong to the \emph{same} connected
component. The sharp contraction and the sharp rebound---symmetric in
magnitude but opposite in sign---are topologically classified together,
not separately.

This is a nontrivial finding and is not what a naive intuition would
predict. A descriptive analysis would tend to separate the contraction
from the rebound, because they have opposite signs and opposite economic
interpretations. The Nada topology, however, treats them as a single
topological unit. The mechanism is the intersection closure: both Q2 and
Q3 2020 appear as members of many of the same base elements of
$\base^{+}$ and $\base^{-}$, because their closed neighborhoods under
the visibility relation overlap substantially in spite of their extreme
and opposite values. Topologically, the pair is a single coherent
episode---the COVID shock as a whole---rather than two unrelated events.
This matches the economic interpretation most commonly given in the
business-cycle literature: the contraction and the rebound of 2020 are
two phases of a single exogenous shock, driven by the same underlying
cause (the policy response to the pandemic), and should be analyzed
jointly rather than as independent events. The topological classifier
recovers this interpretation without any prior economic knowledge,
working only from the geometric structure of the series.

\begin{remark}[What the topology does and does not claim]%
\label{rem:topology-claims}
The finding that observations $114$ and $115$ lie in the same connected
component of the Nada topology is a purely topological statement: it
says that no clopen set of $\tau_{\text{Nada}}$ separates them. It does
\emph{not} say that they are statistically similar in any conventional
sense, nor that their values are close (in fact, they are the two most
distant observations in the value dimension). The topology is not a
proximity measure; it is a structural classifier that responds to the
relational pattern of neighborhoods, not to Euclidean distance. This is
precisely the kind of information that traditional statistical methods
based on correlation, variance, or spectral decomposition cannot
extract, because it lives in the combinatorics of the visibility
relation rather than in the metric structure of the values. The
topological analysis therefore complements, rather than replaces,
standard statistical methods.
\end{remark}

%-------------------------------------------------------------------------------
\subsection{Synthesis: three layers of structure}%
\label{subsec:synthesis}
%-------------------------------------------------------------------------------

The empirical analysis of U.S.\ GDP growth reveals, via the directed
Nada bitopological construction, three layers of structural information
that are individually distinct and jointly informative:

\begin{enumerate}[label=(\arabic*), leftmargin=2em]
\item \textbf{Reachability layer} ($\tau_A$). The DAG is globally
connected: every observation is linked to every later observation
through the compositional structure of the business cycle. One
component, trivial topology at the reachability level.

\item \textbf{Intersection-closure layer} ($\tau_{\text{Nada}}$). Local
fragmentation emerges: the series decomposes into $6$ to $36$ components
depending on graph construction, revealing clusters of observations
that share neighborhood structure beyond mere reachability. The
disconnection is substantive and empirically robust.

\item \textbf{Bitopological layer} $(\taufwd, \taubwd)$. Temporal
asymmetry emerges: forward and backward Nada topologies differ in
component count by $\Delta = +4$ in both HVG and NVG. The forward
topology is more connected than the backward topology, consistent with
the slow-expansion / abrupt-contraction regime of the U.S.\ economy
over this window. The asymmetry is a topological signature of temporal
irreversibility and is invariant to the choice between HVG and NVG.
\end{enumerate}

Each layer is obtained by a different construction---reachability
propagation, intersection closure, and dual intersection closure
respectively---and each layer reveals information not visible at the
other two. The combination of all three constitutes the empirical
content of the directed Nada bitopological analysis of the series.
In this sense, the finding on the COVID-19 shock
(\Cref{subsec:covid-finding}) illustrates what the second layer can do;
the nonzero $\Delta$ values illustrate what the third layer can do; and
the global connectedness of $\tau_A$ sets the baseline of what the first
layer alone would report in isolation.

The entire analysis requires no tuning: the visibility-graph
constructions are parameter-free, the topology-generation procedure is
determined by the combinatorics of the graph, and the specialization
preorder yields exact connected components with no heuristic choices.
The only ``parameter'' is the choice of HVG versus NVG, and for the
purpose of the qualitative findings reported here the two give
consistent answers.


% ══════════════════════════════════════════════════════════════════════════════
\section{Discussion}\label{sec:discussion}
% ══════════════════════════════════════════════════════════════════════════════

The directed extension of the Nada construction developed in this paper
sits at the intersection of several independent research streams: finite
topology (Nada et al.~\cite{Nada2018}, McCord~\cite{McCord1966},
Stong~\cite{Stong1966}), bitopology (Kelly~\cite{Kelly1963}), visibility
graphs for time series (Lacasa et al.~\cite{Lacasa2008},
Luque et al.~\cite{Luque2009}), and the statistical analysis of temporal
irreversibility (Lacasa and Toral~\cite{LacasaToral2010}). We now place
our contribution explicitly against these streams, discuss the epistemic
status of the formal proofs and their empirical realization, acknowledge
the limitations of the framework, and outline directions for further
work.

%-------------------------------------------------------------------------------
\subsection{Relation to prior work}\label{subsec:relation-prior}
%-------------------------------------------------------------------------------

\paragraph{Nada, El Atik and Atef (2018).}
The undirected Nada construction~\cite{Nada2018} was introduced as a
systematic way to associate topological structures with arbitrary
undirected graphs, by taking closed neighborhoods as a subbase and
closing under finite intersections and arbitrary unions. The
contribution of the original paper was the identification of the
family $\{\Ncl{v} : v \in V\}$ as a topologically meaningful generator,
and the study of the resulting topological invariants---connectedness,
separation, continuity between graph morphisms---as graph properties.
Our contribution with respect to Nada et al.\ is twofold. First, we
observe that the construction is \emph{indifferent} to whether the
generating family consists of undirected or directed neighborhoods,
because the finite-intersection and arbitrary-union operations used in
the construction are universal operations on the powerset and do not
depend on any symmetry property of the underlying graph
(Theorem~\ref{thm:directed-validity}). This observation is, as we have
emphasized, Before elaborating, we state the rhetorical hierarchy plainly. The proof is a one-line corollary of the universality of the generated-topology functor on $\mathcal{P}(V)$ and carries essentially no mathematical content; the contribution of this paper begins, rather than ends, at this corollary. structurally trivial once stated correctly: it is a
one-sentence corollary of the universality of the generated-topology
functor on $\powerset(V)$. The triviality is the point: the extension
is available \emph{for free}, with no new mathematical machinery. What
is not trivial is the second component of our contribution: the
recognition that when the construction is applied separately to the
forward and backward neighborhood families of a directed graph, the
pair $(\taufwd, \taubwd)$ is not merely two topologies but a
bitopological space in the sense of Kelly, and that its pairwise
properties encode temporal asymmetry.

\paragraph{Kelly (1963).}
Kelly's theory of bitopological spaces~\cite{Kelly1963} was developed
in a purely formal context, with no particular empirical application
in mind. The natural examples at the time of its introduction were
metric-theoretic---quasi-pseudometric spaces and their forward and
backward topologies---and algebraic. To our knowledge, the
bitopological structure has not previously been used to analyze time
series or to quantify temporal irreversibility. Our contribution is to
\emph{realize} Kelly's abstract bitopology on a concrete dynamical
object: the directed visibility graph of a time series. The forward
topology $\taufwd$ arises from forward neighborhoods that look toward
the future; the backward topology $\taubwd$ arises from backward
neighborhoods that look toward the past. The pair
$(V, \taufwd, \taubwd)$ is an honest bitopological space, and its
pairwise connectedness characterizes whether the forward and backward
views of the series are topologically compatible. The irreversibility
indices of Proposition~\ref{prop:irreversibility-indices} quantify the
degree to which they are not.

\paragraph{Lacasa et al.\ (2008) and Luque et al.\ (2009).}
Visibility graphs map time series into combinatorial objects that
preserve dynamical structure: periodic series yield regular graphs,
random series yield exponential degree distributions, fractal series
yield power-law degree distributions. The directed versions of these
graphs preserve the temporal orientation of the series, so that causal
precedence is retained. Prior work on directed visibility graphs has
focused primarily on degree-based asymmetry statistics
(in-degree versus out-degree distributions, directed clustering) as a
measure of temporal irreversibility. Our contribution is orthogonal to
those metrics: the bitopological analysis operates at the level of the
full neighborhood structure, not just the degree counts, and captures
asymmetries that are invisible to degree statistics. In particular, as
our empirical analysis of the U.S.\ GDP series demonstrates
(\Cref{sec:empirical}), the base sizes $|\base^{+}|$ and $|\base^{-}|$
can be exactly equal---yielding $I_{B} = 0$---while the component
structures of $\taufwd$ and $\taubwd$ differ. This is an asymmetry that
lives in the combinatorics of the intersection closure, not in the
cardinality of the generators, and it would be missed by any method
that examines only the first-order neighborhood statistics.

\paragraph{McCord (1966) and Stong (1966).}
The specialization-preorder theorem for finite topological
spaces~\cite{McCord1966,Stong1966} is the technical instrument that
makes our framework computationally tractable.
\Cref{thm:specialization-components} reduces topological connectedness---a
property defined in terms of all open sets---to connectedness of an
auxiliary graph built entirely from the base. The consequence is that
decision of connectedness becomes polynomial in $n$ and $|\base|$,
with no dependence on the potentially exponential size of the
topology. Without this reduction, the directed Nada bitopological
analysis would be computationally infeasible for anything beyond trivial
series lengths. The theorem is classical but, as we noted in the Lean
formalization, it is not currently available in Mathlib; our single
remaining \texttt{sorry} in the formalization corresponds exactly to
this result, and closing that gap---either by formalizing McCord's
original proof in Lean 4 or by deriving it from the existing Mathlib
machinery for specialization preorders---is a natural piece of future
work.

\paragraph{Lacasa and Toral (2010).}
The analysis of temporal irreversibility in time series via directed
visibility graphs was pioneered by Lacasa and Toral~\cite{LacasaToral2010},
who proposed measuring irreversibility through the asymmetry of the
degree distributions of the in- and out-neighborhoods. Their framework
is information-theoretic: the irreversibility of the series is
quantified by the Kullback-Leibler divergence between the in-degree
and out-degree distributions, and the resulting index is zero for
reversible dynamics and strictly positive for irreversible ones. Our
bitopological framework yields a parallel quantification: the
component-count irreversibility index~$I_C$ and the asymmetry
direction~$\Delta$ of
Proposition~\ref{prop:irreversibility-indices} are topological
analogues of their statistical asymmetry measures. The two approaches
are complementary: Lacasa and Toral measure irreversibility at the
level of neighborhood size statistics, we measure it at the level of
topological component structure. On the U.S.\ GDP series our
bitopological framework yields $\Delta = +4$ in both the HVG and NVG
constructions, a value that is stable across graph types and points in
the direction predicted by the Lacasa-Toral theory for series with
slow expansions and abrupt contractions.

%-------------------------------------------------------------------------------
\subsection{Formalist compression and constructive content}%
\label{subsec:formalism-constructivism}
%-------------------------------------------------------------------------------

A methodological remark deserves emphasis. The validity of the directed
extension of the Nada construction is, as we have seen, provable in a
handful of lines: the topology axioms are universal operations on the
powerset $\powerset(V)$, and the operations of closure under finite
intersection and arbitrary union do not inspect whether the generating
family was produced by a symmetric, antisymmetric, transitive, or
arbitrary graph
(Theorem~\ref{thm:directed-validity},
Theorem~\ref{thm:bitopological-structure}). Once one has identified the
correct structural abstraction---namely, that the Nada construction is
a pair of instances of the functor $\mathcal{S} \mapsto
\mathrm{generateFrom}(\mathcal{S})$ from families of subsets to
topologies---the directed extension is a corollary of two applications
of the same functor, one to $\{\Nfwd{v}\}$ and one to $\{\Nbwd{v}\}$.
In our Lean 4 formalization against Mathlib, the proof reduces
literally to two invocations of
\texttt{TopologicalSpace.generateFrom}, each automatically yielding a
topology by construction.

This compressed, formalist proof is entirely satisfactory as a
\emph{justification of validity}. It establishes, once and for all,
that whatever directed graph we feed into the procedure, we obtain two
topologies on its vertex set, and that their pair is a bitopological
space. No symmetry hypothesis is required; no graph-theoretic property
is assumed; the proof applies uniformly to every directed graph.

What the compressed proof does \emph{not} tell us is whether the
construction reveals anything substantive about the data on which it is
run. The proof is indifferent to the content of the series. The
empirical question---whether the bitopological asymmetry captures
meaningful features of the dynamics, whether the component structure
distinguishes economically interpretable episodes, whether the
resolution gain between the Alexandrov and Nada layers is genuinely
informative or merely notational---cannot be answered from the
formalist proof alone. It can be answered only by running the algorithm
on real data and observing what it returns.

The algorithms of \Cref{sec:algorithms} are, in this sense, the
\emph{decompressed realization} of the formalist proof. They take the
abstract claim ``the Nada construction extends to directed graphs and
produces a bitopological space'' and convert it into a procedure that,
given a specific time series of length $n$, returns specific component
counts, specific base sizes, specific irreversibility indices, and a
specific answer to whether observations $114$ and $115$ of the U.S.\
GDP series lie in the same connected component. The polynomial-time
decidability of connectedness via the specialization preorder
(\Cref{thm:specialization-components}) is what makes this decompression feasible;
without it, the formalist guarantee would be of purely theoretical
interest.

The relation between the formalist proof and the constructive algorithm
is therefore not one of opposition or rivalry. It is one of two
complementary moments of a single methodological process: the proof
establishes that the construction is available for every input, and the
algorithm supplies what the proof, by its nature, cannot supply---the
actual answer on any specific input. The validity of the method
depends on the proof; the \emph{content} of the method depends on the
algorithm. Neither is dispensable, and neither subsumes the other.

This reading has a further consequence. The three-valued output of
\Cref{alg:pairwise}---returning \textsc{connected},
\textsc{disconnected}, or \textsc{undecided} when the resource budget is
exhausted---is sometimes criticized as an inelegance: why not default
to one of the two Boolean outcomes when the enumeration cannot complete?
\Cref{subsec:pairwise-algorithm} in \Cref{sec:algorithms} gave the
Brouwer--Heyting--Kolmogorov answer: because on a finite space the
proposition is decidable \emph{in principle} but its decision
procedure is a concrete computation, and that computation either has
been run to completion or it has not. When it has not, claiming either
outcome is an extrapolation that the evidence does not support. The
honest report is the epistemic state of the algorithm after finite
work, not a guess about what it would have returned if it had been
allowed to run longer. This is the discipline of treating propositions
as computational tasks: the proof \emph{is} the completed task, and
when the task is incomplete there is no proof yet. The algorithm
reports this faithfully, and the user is free to allocate more
resources if a definite answer is required.

%-------------------------------------------------------------------------------
\subsection{Limitations}\label{subsec:limitations}
%-------------------------------------------------------------------------------

Several limitations of the present work should be stated explicitly.

\paragraph{Pairwise connectedness requires topology enumeration.}
The decision of pairwise connectedness
(\Cref{alg:pairwise}) operates on the fully enumerated
topologies $\taufwd$ and $\taubwd$. Unlike connectedness of a single
topology, which is decidable from the base alone via the specialization
preorder, pairwise connectedness has no analogous polynomial-time
shortcut that we are aware of: the decision genuinely requires
searching among the open sets for a disconnecting clopen, and the
number of open sets can grow as $2^{|\base|}$ in the worst case. For
the time series analyzed in \Cref{sec:empirical}, with base sizes
around $250$, this places full enumeration well beyond feasible. In
practice, the algorithm is invoked with a safety bound
$\texttt{max\_open\_sets}$ and, when the bound is exceeded, honestly
returns the undecided outcome. Developing a tractable decision
procedure for pairwise connectedness that works directly from the two
bases, without topology enumeration, would eliminate this limitation
and is an open problem. We conjecture that an analogue of
\Cref{thm:specialization-components} for the bitopological specialization
structure may exist and would yield such a procedure, but we have not
proved this.

\paragraph{The resolution gain is empirical, not theoretical.}
Throughout the paper, we have emphasized that the positive gap
$|\base_{\text{Nada}}| - |\base_{A}|$ observed on the empirical series
is a fact about the specific generators used in the two pipelines and
not a corollary of $\tau_A \subseteq \tau^{+}_{\text{Nada}}$
(Theorem~\ref{thm:alexandrov-subtopology}). The inclusion of
\emph{topologies} is a theorem; the inequality of \emph{base sizes} is
an observation that holds on the empirical series we analyzed and need
not hold universally. On other inputs the direction of the gap could in
principle be reversed. We do not know a structural characterization of
the class of directed graphs on which the Nada intersection closure
produces strictly more base elements than the Alexandrov reachability
pipeline, and providing such a characterization is another open
question.

\paragraph{A classical result is not yet formalized in Mathlib.}
Our Lean 4 formalization of Sections~\ref{sec:prelim} through
\ref{sec:extension} compiles cleanly against the current Mathlib, with
exactly one \texttt{sorry}: the statement of
\Cref{thm:specialization-components}, the McCord--Stong equivalence between
topological connectedness and comparability-graph connectedness for
the specialization preorder of a finite space. The theorem is
classical, but its formalization against Mathlib requires the
minimal-open-set characterization of points in a finite Alexandrov
space and the equivalence between topological and comparability-graph
connectedness in that setting, neither of which is currently in
Mathlib's library. We have left the theorem as a stated but unproved
lemma (with a comment pointing to the classical references), and
every theorem of the formalization that depends on it carries the
dependency explicitly. The gap is purely a matter of Mathlib coverage,
not of mathematical uncertainty, and we expect that a dedicated effort
of a few hundred lines of Lean code would close it. We would welcome
a Mathlib contribution implementing the McCord--Stong theorem
systematically, as it has applications well beyond the present work.

\paragraph{The method is developed for finite ground sets.}
The specialization preorder is defined on any topological space, and
the Alexandrov correspondence between finite preorders and finite
Alexandrov topologies is a special case of a more general adjunction.
In principle the directed Nada construction extends to any space on
which closed neighborhoods are defined. In practice the base closure
under finite intersection may fail to terminate in the infinite case,
or the resulting base may fail to be a topological base in the
standard sense. A systematic treatment of the infinite case would
require additional hypotheses---local finiteness of the graph,
compactness of the ground space, or similar---and we have not
attempted it here. For time-series analysis the finite case is
sufficient, since every observed series has finite length.

%-------------------------------------------------------------------------------
\subsection{Extensions and future work}\label{subsec:future-work}
%-------------------------------------------------------------------------------

The directed Nada bitopological analysis opens several directions for
continued research.

\paragraph{Directed graphs beyond visibility graphs.}
The construction of \Cref{sec:extension} applies to any directed
graph, and nothing in the proof of validity depends on the graph
arising from a visibility relation. Natural candidates for further
application include: causal graphs in epidemiology and econometrics,
where the directed edges represent temporal precedence of events;
dependency graphs in software engineering and project management,
where edges represent prerequisite relations; transaction networks in
finance and blockchain analysis, where edges represent directional
flows of value; and gene regulatory networks in systems biology, where
edges represent inferred activation or repression relations. In each
of these settings, the bitopological analysis would yield forward and
backward topologies whose asymmetry would quantify a domain-specific
notion of irreversibility or directional structure. Whether the
resulting invariants are useful for the substantive questions of each
domain is an empirical question that can only be answered by
application.

\paragraph{Weighted graphs and filtrations.}
The construction of this paper treats directed edges as unweighted.
A natural generalization is to consider a weighted directed graph
together with a threshold $\theta$, construct the unweighted graph
obtained by retaining only edges of weight at least $\theta$, and
apply the Nada construction. Varying $\theta$ yields a
\emph{filtration} of bitopological spaces, indexed by the threshold.
The evolution of the invariants
$(C^{+}(\theta), C^{-}(\theta), \Delta(\theta))$ as $\theta$ varies
would yield a persistent-bitopology analogue of persistent homology.
The connection to persistent homology is potentially deep and warrants
systematic investigation, particularly because the bitopological
setting introduces a direction-sensitive refinement not present in the
standard (undirected, scalar-valued) persistent-homology framework.

\paragraph{Tighter formalization in Mathlib.}
As noted above, closing the single \texttt{sorry} in our Lean
formalization would complete the mechanical verification of all
theorems stated in this paper. Beyond this, a systematic development
of finite-space bitopology in Mathlib---including the specialization
preorder for bitopological spaces, the pairwise variants of separation
and connectedness, and the associated algorithms---would provide
infrastructure for further formalizations of time-series topology and
would establish a pattern for integrating combinatorial and continuous
topology in a single library.

\paragraph{Applications to other time series with known asymmetry.}
The U.S.\ GDP series analyzed in \Cref{sec:empirical} was selected as
a paradigmatic example of asymmetric economic dynamics. Other series
with known temporal asymmetries would provide further tests of the
method: the daily closing prices of equity indices during regimes of
flash crashes and gradual recoveries; seismic signals during the
buildup and release of fault stress; turbulent fluid velocity
fluctuations with intermittent burst events; and neuronal spike
trains during transitions between quiescent and active states.
Systematic application of the bitopological framework to a diverse
portfolio of such series, together with comparison against the
degree-based asymmetry statistics of Lacasa and
Toral~\cite{LacasaToral2010}, would clarify the relative strengths of
the topological and statistical approaches and may suggest hybrid
methods that combine their respective advantages.

% ══════════════════════════════════════════════════════════════════════════════
\section{Conclusion}\label{sec:conclusion}
% ══════════════════════════════════════════════════════════════════════════════

We have shown that the topological construction of
Nada et al.~\cite{Nada2018}, originally formulated for undirected
graphs, extends verbatim to directed graphs and produces a
bitopological space in the sense of Kelly~\cite{Kelly1963} when
applied separately to the forward and backward closed neighborhoods.
The extension is formally trivial---a two-line corollary of the
universality of the generated-topology functor on the powerset---but
it opens a substantive new framework for the analysis of temporally
asymmetric dynamical systems. Applied to directed visibility graphs of
time series, the construction yields a hierarchy of three topological
layers: the Alexandrov topology of the reachability preorder, which
captures global connectedness at the level of causal precedence; the
Nada topology of the undirected closed neighborhoods, which reveals
local fragmentation invisible at the reachability layer; and the
directed Nada bitopological pair, whose asymmetry directly quantifies
temporal irreversibility. We have proved the validity of the extension
and a family of related results on the Alexandrov sublayer, the
specialization preorder, the three irreversibility indices, and
pairwise connectedness; formalized the proofs in Lean~4 against
\texttt{mathlib4}~\cite{mathlib} without any \texttt{sorry}
placeholders, relying on exactly one explicitly declared axiom
corresponding to the finite-space form of the classical McCord--Stong
theorem~\cite{McCord1966,Stong1966} (not yet available in
\texttt{mathlib4}, and explicitly discoverable via
\texttt{\#print axioms}); supplied a complete suite of algorithms with
polynomial-time complexity analysis and a three-valued decision
procedure for pairwise connectedness that reports honestly when its
resource budget is exhausted; and demonstrated the empirical content
of the framework on the quarterly U.S.\ real GDP series over
1992--2024, recovering a positive asymmetry direction consistent with
the slow-expansion / abrupt-contraction regime of the economy and a
coherent topological classification of the COVID-19 contraction and
rebound as a single connected episode. The formalist proof and the
constructive algorithm are two moments of the same methodological
process: the proof establishes availability, the algorithm supplies
content, and the two together give a framework that is both rigorously
justified and empirically meaningful.



% ══════════════════════════════════════════════════════════════════════════════
\section*{Data and Code Availability}
\addcontentsline{toc}{section}{Data and Code Availability}
% ══════════════════════════════════════════════════════════════════════════════

All code, data, and formal proofs supporting the results of this paper
are archived at the public repository
\url{https://github.com/IsadoreNabi/bitopology-nada} and permanently
deposited on Zenodo~\cite{bitopology-nada-zenodo} under the concept DOI
\href{https://doi.org/10.5281/zenodo.19488414}{\texttt{10.5281/zenodo.19488414}},
which resolves to the latest version; the specific version used in
this paper is v1.0.0
(\href{https://doi.org/10.5281/zenodo.19488415}{\texttt{10.5281/zenodo.19488415}}).
The repository contains: (i) the \LaTeX{} source of this paper; (ii) an
R~Markdown document \texttt{reproducibility.Rmd} that regenerates
Table~\ref{tab:gdp-results} from the FRED series \texttt{GDPC1}~\cite{FRED-GDPC1}
using the \texttt{topologyR} package~\cite{topologyR}, together with a
frozen CSV snapshot of the exact data vintage used here so that
results are bit-exact reproducible under later BEA revisions;
(iii) the Lean~4 formalization of
Sections~\ref{sec:prelim}--\ref{sec:extension} against
\texttt{mathlib4}~\cite{mathlib}, which contains no \texttt{sorry}
placeholders and relies on exactly one explicitly declared axiom
(\texttt{specialization\_chain\_of\_sameComponent}) corresponding to the
finite-space form of the McCord--Stong
theorem~\cite{McCord1966,Stong1966}, classical but not yet available in
\texttt{mathlib4}; this dependency is explicitly discoverable via
Lean's \texttt{\#print axioms} command on any downstream theorem; and
(iv) a README describing the three independent compilation workflows
(\texttt{pdflatex}, \texttt{Rscript}, \texttt{lake build}). The
empirical results of Section~\ref{sec:empirical}, including the
$24{+}6$ soft checks documented in the R~Markdown source, are
verifiable by a single command.


\begin{thebibliography}{99}

\bibitem{Alexandrov1937}
P.~S.~Alexandrov,
\emph{Diskrete R\"aume},
Matematicheskii Sbornik \textbf{2}(44), no.~3, 501--519, 1937.

\bibitem{Kelly1963}
J.~C.~Kelly,
\emph{Bitopological spaces},
Proceedings of the London Mathematical Society, Series~3, \textbf{13},
71--89, 1963.

\bibitem{Lacasa2008}
L.~Lacasa, B.~Luque, F.~Ballesteros, J.~Luque, and J.~C.~Nu\~no,
\emph{From time series to complex networks: the visibility graph},
Proceedings of the National Academy of Sciences \textbf{105}(13),
4972--4975, 2008.

\bibitem{LacasaToral2010}
L.~Lacasa and R.~Toral,
\emph{Description of stochastic and chaotic series using visibility graphs},
Physical Review E \textbf{82}(3), 036120, 2010.

\bibitem{Luque2009}
B.~Luque, L.~Lacasa, F.~Ballesteros, and J.~Luque,
\emph{Horizontal visibility graphs: exact results for random time series},
Physical Review E \textbf{80}(4), 046103, 2009.

\bibitem{McCord1966}
M.~C.~McCord,
\emph{Singular homology groups and homotopy groups of finite topological spaces},
Duke Mathematical Journal \textbf{33}(3), 465--474, 1966.

\bibitem{Nada2018}
S.~Nada, A.~E.~F.~El~Atik, and M.~Atef,
\emph{New types of topological structures via graphs},
Mathematical Methods in the Applied Sciences \textbf{41}(15), 5801--5810,
2018.

\bibitem{Stong1966}
R.~E.~Stong,
\emph{Finite topological spaces},
Transactions of the American Mathematical Society \textbf{123}(2),
325--340, 1966.

\bibitem{topologyR}
J.~M.~Gomez Julian,
\emph{\texttt{topologyR}: Topology Generation from Graphs in R},
R package, 2026.
\url{https://github.com/IsadoreNabi/topologyR}.

\bibitem{mathlib}
The mathlib Community,
\emph{The Lean Mathematical Library},
in Proceedings of the 9th ACM SIGPLAN International Conference on
Certified Programs and Proofs (CPP 2020), 367--381, 2020.
\url{https://doi.org/10.1145/3372885.3373824}.

\bibitem{FRED-GDPC1}
U.S.~Bureau of Economic Analysis,
\emph{Real Gross Domestic Product} [GDPC1],
retrieved from FRED, Federal Reserve Bank of St.~Louis,
\url{https://fred.stlouisfed.org/series/GDPC1}, accessed April 2026.

\bibitem{bitopology-nada-zenodo}
J.~M.~Gomez Julian,
\emph{Bitopological Spaces from Directed Graphs: Reproducibility and
Formalization Bundle}, version 1.0.0, Zenodo, 2026.
\url{https://doi.org/10.5281/zenodo.19488414}.

\end{thebibliography}

\end{document}
